\documentclass[11pt]{ctexrep}

% ---------------------------------------------------------------------------
% normet 用户手册（中文版）——直接用中文写成，章节、标签与 normet_manual.tex 一一对应。
% 构建：docs/manual/build.sh（XeLaTeX + ctex）
% 截图：docs/manual/make_figures.py，NORMET_UI_LANGUAGE=zh 时存入 figures_zh/；
%       没有中文截图的取 figures/。
% ---------------------------------------------------------------------------
\usepackage[a4paper,margin=1in]{geometry}
\usepackage{amsmath,amssymb}
\usepackage{graphicx}
\usepackage{booktabs}
\usepackage{longtable}
\usepackage{array}
\usepackage{xcolor}
\usepackage{fancyhdr}
\usepackage{listings}
\usepackage[hidelinks,colorlinks=true,linkcolor=blue!55!black,urlcolor=blue!55!black,
            citecolor=blue!55!black]{hyperref}

% Compact list spacing without the enumitem package.
\newcommand{\nosep}{\setlength{\itemsep}{0pt}\setlength{\parskip}{0pt}\setlength{\topsep}{2pt}}

\definecolor{nmblue}{HTML}{2C7BB6}
\definecolor{nmgray}{HTML}{8E8E93}
\definecolor{codebg}{HTML}{F5F5F7}
\definecolor{warn}{HTML}{FF9500}

% --- code listings -----------------------------------------------------------
% A listing's comments may be Chinese: listings passes them through as they are.
\lstdefinestyle{normet}{
  backgroundcolor=\color{codebg},
  basicstyle=\ttfamily\small,
  keywordstyle=\color{nmblue}\bfseries,
  commentstyle=\color{nmgray}\itshape,
  stringstyle=\color{green!45!black},
  breaklines=true, frame=single, rulecolor=\color{nmgray!40},
  showstringspaces=false, upquote=true, columns=fullflexible, xleftmargin=4pt, framexleftmargin=4pt,
  extendedchars=false,
}
\lstset{style=normet}

% --- headers / footers -------------------------------------------------------
\pagestyle{fancy}
\fancyhf{}
\fancyhead[L]{\small normet 用户手册}
\fancyhead[R]{\small\leftmark}
\fancyfoot[C]{\thepage}
\renewcommand{\headrulewidth}{0.4pt}

% --- handy macros ------------------------------------------------------------
\newcommand{\nm}{\textsf{normet}}
\newcommand{\codeunderscore}{\textunderscore\allowbreak}
\newcommand{\code}[1]{\texttt{\let\_\codeunderscore #1}}
\setlength{\emergencystretch}{2.5em}
\newcommand{\tab}[1]{\textbf{#1}}
\newcommand{\NOtwo}{NO\textsubscript{2}}
\newcommand{\PMtwo}{PM\textsubscript{2.5}}
\newcommand{\ugm}{\ensuremath{\mu}g\,m\textsuperscript{$-3$}}
\newcommand{\note}[1]{\par\noindent\colorbox{codebg}{\parbox{0.97\linewidth}{\small\textbf{注意。}#1}}\par}
\newcommand{\caution}[1]{\par\noindent\colorbox{warn!15}{\parbox{0.97\linewidth}{\small\textbf{警告。}#1}}\par}
% The window as the Chinese manual shows it: its screenshot in Chinese where one was
% taken (figures_zh/), else the English one.
\newcommand{\screenshot}[3]{%
  \begin{figure}[htbp]\centering
    \fboxsep=0pt\fbox{\IfFileExists{figures_zh/#1}%
      {\includegraphics[width=0.98\linewidth]{figures_zh/#1}}%
      {\includegraphics[width=0.98\linewidth]{figures/#1}}}%
    \caption{#2}\label{#3}
  \end{figure}}

\begin{document}

% ============================ TITLE PAGE ====================================
\begin{titlepage}
\centering
\vspace*{2.5cm}
{\IfFileExists{../../packaging/assets/normet.png}{\includegraphics[width=3.2cm]{../../packaging/assets/normet.png}\\[1cm]}{}}
{\Huge\bfseries\color{nmblue} normet\par}
\vspace{0.4cm}
{\Large 环境时间序列的归一化、分解\\与反事实建模\par}
\vspace{1.2cm}
{\large\bfseries 用户手册\par}
\vspace{0.3cm}
{\large 版本 1.0.0\par}
\vfill
{\large 一个桌面应用，一个 Python 库：从空气质量记录中剥离天气的影响，\\
并回答一项干预究竟改变了什么。\par}
\vspace{1.5cm}
{\color{nmgray}\today}
\end{titlepage}

\tableofcontents
\cleardoublepage

% ============================ 1. INTRODUCTION ===============================
\chapter{简介}\label{ch:intro}

\section{normet 是什么？}
\nm{} 是一个桌面应用，也是一个 Python 库，用来分析环境时间序列——首先是空气质量记录。
原始浓度记录回答不了的两个问题，它都回答得了：
\begin{enumerate}\nosep
  \item \emph{我看到的这些变化，有多少是天气造成的？}一个月的静稳、寒冷反气旋天气，
        和新添一处污染源一样，都能把浓度抬上去。\nm{} 先学习污染物如何依赖气象与日历，
        再把气象平均掉，留下天气解释不了的那部分记录。这就是\emph{气象归一化}，
        也叫\emph{去气象化}。
  \item \emph{假如没有这项干预，会是什么样？}低排放区、封控或工厂关停，某一天在某个
        地方改变了局面。\nm{} 拿干预没有触及的其他地方，为这个地方构造一个反事实——
        \emph{合成控制}——再度量干预日期之后的差异，并用安慰剂检验和不确定性区间判断
        这个差异是否超出偶然。
\end{enumerate}
窗口里的每一个操作，背后都是对 \nm{} Python 包的一次调用；这些调用可以写成脚本、在命令行
运行，也可以经由模型上下文协议（MCP）交给 AI 助手驱动（第\ref{ch:script}章）。窗口会记下
自己做过的每一次调用，所以点出来的分析也能存成一个脚本，原样重跑
（第~\ref{sec:commandlog}~节）。

\section{一个方程讲清思路}
设 $y_t$ 为 $t$ 时刻的浓度，$\mathbf{m}_t$ 为该时刻的气象（风速风向、温度、湿度、边界层
高度等），$\mathbf{c}_t$ 为日历（趋势计数、一年中的第几天、星期几、小时）。\nm{} 先拟合一个
机器学习模型
\begin{equation}
  y_t \approx f(\mathbf{m}_t, \mathbf{c}_t)
\end{equation}
然后追问：若把 $t$ 时刻的天气换成该站点曾经出现过的任意一种天气，$f$ 会预测出什么？把许多
次这样的抽样平均起来，
\begin{equation}
  \tilde{y}_t \;=\; \frac{1}{N}\sum_{k=1}^{N} f\bigl(\mathbf{m}^{(k)}, \mathbf{c}_t\bigr),
  \qquad \mathbf{m}^{(k)} \sim \text{记录自身的气象},
  \label{eq:norm}
\end{equation}
就得到\emph{归一化}序列 $\tilde{y}_t$，即典型天气下 $t$ 时刻应有的浓度。日历 $\mathbf{c}_t$
保留不动，排放的变化——模型只能通过趋势、季节、星期和小时来表达它——便留在 $\tilde{y}_t$
里；天气造成的逐日起伏则被平均掉。第\ref{ch:theory}章展开这条思路，以及建立在它之上的各种
分解。

\section{主要功能}
\begin{itemize}\nosep
  \item AutoML 训练：FLAML 在墙钟\emph{预算}内搜索估计器，或使用调参的 LightGBM；以
        $R^2$、RMSE 及另外九项统计量报告留出集上的拟合。
  \item 蒙特卡洛气象归一化：重采样哪些变量可选，另有不确定性分位数、参考池和条件
        重采样；自适应版本会在序列收敛后自动停止。
  \item 排放分解（趋势、季节、星期、日变化、短期）和气象分解（每个变量或命名分组一项
        贡献，顺序归因或 Shapley 归因）。
  \item 滚动（分窗）归一化，以及按时间尺度区分天气影响的多尺度分解。
  \item 已训练模型的部分依赖曲线。
  \item 用 Chronos-2 时间序列基础模型做零样本去气象化，结果格式与上面相同——不必训练。
  \item 面板数据上的合成控制：增广 SCM、经典 Abadie SCM、双重差分、矩阵补全、稳健 SCM、
        ML-SCM 和贝叶斯 SCM，配有空间与时间安慰剂检验、刀切法和自助法区间，并可对全部
        单元批量运行。
  \item 数据工作室：来自六个英国网络、覆盖全欧的 EEA 报告网络或覆盖世界大部分地区的
        OpenAQ 的空气质量数据，加再分析气象，合并成可建模的表（英国网络和 EEA 无需
        API 密钥）。
  \item 传输工作室：把 HYSPLIT 后向轨迹转换成传输预测变量，合并到记录上。
  \item 命令日志：每个操作都写成能重复它的调用；带溯源信息的 HTML 报告；按期刊栏宽
        排版的图。
  \item 助手：一个在窗口里工作的语言模型，提议目标变量和预测变量，陪你解读结果。
\end{itemize}

\section{本手册的组织}
第\ref{ch:install}章讲安装。第\ref{ch:theory}章讲方法，深到足以解读程序报出的每一个数字。
第\ref{ch:start}--\ref{ch:pdp}章在窗口里用一条真实记录走完一次完整的去气象化分析，
第\ref{ch:scm}章再用一个真实面板走完合成控制流程。第\ref{ch:data}章讲数据导入，
第\ref{ch:export}章讲导出、报告和重现一次会话，第\ref{ch:assistant}章讲助手，
第\ref{ch:script}章讲脚本。附录列出默认值、\nm{} 存放文件的位置，以及故障排除指南。

书中的图都是真实窗口的截图，由 \code{docs/manual/make\_figures.py} 在 \nm{} 自带的记录上
生成：伦敦 Marylebone Road（站点 MY1）2020 年 1 月至 8 月的逐时 \NOtwo{} 和 ERA5 气象，
这段时间跨越了英国第一次 COVID-19 封控。正文引用的数字就是该脚本打印出来的数字。你自己
跑出来的结果会在最后一两位上不同：AutoML 搜索以墙钟时间为界，两次访问的候选模型不会完全
一样。

本手册另有英文版（\code{normet\_manual.pdf}），章节一一对应，截图为英文界面
（第~\ref{sec:appearance}~节）。

% ============================ 2. INSTALLATION ===============================
\chapter{安装}\label{ch:install}

\nm{} 为每个平台提供一个自包含的桌面应用，PyPI 上另有 Python 包。做交互式分析，装应用；
写脚本、上服务器、用笔记本，装 Python 包。

\section{macOS}
打开 \code{Normet-1.0.0-macos-arm64.dmg}，把 \code{Normet.app} 拖进\code{应用程序}文件夹。
应用自带 Python、Qt、FLAML、LightGBM 和 XGBoost，以及数据工作室要用的读取器，所以去气象化、
合成控制和数据获取都不用另装东西。首次启动比以后慢，因为要建字体缓存；之后的启动复用它
（保存在 \code{\textasciitilde/.normet/matplotlib}）。建立期间会显示启动画面。

\note{应用没有用 Apple 开发者 ID 签名。第一次打开时，右键点击 \code{Normet.app}，选
\emph{打开}，在对话框里确认一次即可。也可以在终端运行
\code{xattr -dr com.apple.quarantine /Applications/Normet.app}。}

安装包不含可选的 Chronos-2 基础模型（它依赖 PyTorch，光这一项就约 1~GB）、贝叶斯 SCM 后端
（PyMC）和 MCP 服务器。窗口里它们的选项呈灰色，工具提示会说明缺什么；要用，请从 PyPI 安装
带相应扩展的 \nm{}（见下文）。

\section{Windows 和 Linux}
Windows 上运行 \code{normet-setup-1.0.0.exe}。这是一个 Inno~Setup 安装向导，会添加开始菜单
快捷方式和可选的桌面快捷方式。安装程序未签名，首次运行时 SmartScreen 会显示\emph{Windows 已
保护你的电脑}：选\emph{更多信息} $\rightarrow$ \emph{仍要运行}。Linux 上下载
\code{normet-1.0.0-x86\_64.AppImage}，赋予可执行权限（\code{chmod +x}）后运行。精简的发行版
上，Qt 可能需要 \code{libegl1 libgl1 libxkbcommon0 libdbus-1-3 libxcb-cursor0}。

安装包由仓库的 \code{build-gui.yml} 工作流构建（每个操作系统一个；PyInstaller 不能交叉编译），
附在各次运行记录的 \emph{Actions $\rightarrow$ Artifacts} 下。\code{packaging/README.md}
说明了如何在本地构建。

\section{从 PyPI 安装}
\nm{} 需要 Python~3.10 或更高版本（测试过 3.10--3.12），NumPy 须低于 2。核心包很小，较重的
依赖都做成扩展：
\begin{lstlisting}[language=bash]
python -m venv .venv && source .venv/bin/activate
pip install "normet[gui,flaml]"      # the desktop app with the AutoML backend
normet-gui                           # or: normet-gui mydata.csv
\end{lstlisting}

\begin{center}\small
\begin{tabular}{@{}ll@{}}
\toprule
扩展 & 增加的功能 \\
\midrule
\code{flaml}      & FLAML AutoML 后端（含 FLAML 会导入的 XGBoost）——推荐 \\
\code{lgb}        & LightGBM 后端，随机超参数搜索 \\
\code{gui}        & PySide6，用于 \code{normet-gui} \\
\code{data}       & ERA5（CDS）、英国监测网络、OpenAQ、EEA、Open-Meteo 适配器 \\
\code{geo}        & Shapely 和 pyshp：从 GeoJSON/Shapefile 读取源区 \\
\code{cli}        & Click 和 PyYAML，用于 \code{normet} 命令 \\
\code{foundation} & Chronos-2 和 PyTorch：零样本去气象化、站点嵌入 \\
\code{finetune}   & \code{foundation} 加 \code{accelerate} 和 \code{peft}，用于微调 Chronos-2 \\
\code{bayes}      & PyMC 和 ArviZ，用于贝叶斯 SCM \\
\code{mcp}        & MCP 服务器 \code{normet-mcp}，供 AI 助手使用 \\
\code{dask}       & Dask，用于大数据集 \\
\code{all}        & 以上全部 \\
\bottomrule
\end{tabular}
\end{center}

导入 \nm{} 从不要求可选依赖：某个函数需要而没装时，它会告诉你装哪个扩展。

% ============================ 3. METHODOLOGY =================================
\chapter{方法}\label{ch:theory}

本章说明 \nm{} 到底算了什么，好让你读它报出的数字时心里有底。跑一次分析不必读；要替一次
分析辩护，就得读。

\section{准备数据}\label{sec:prepare}
训练之前，\code{prepare\_data} 按下面的顺序处理表格：
\begin{enumerate}\nosep
  \item \textbf{日期列。}表里必须有一列日期时间，名为 \code{date}（也可以是索引，或名为
        \code{datetime}、\code{time}、\code{timestamp} 的列，会被改名）。带时区的日期先
        换算成 UTC，再去掉时区。
  \item \textbf{目标和预测变量。}目标改名为 \code{value}；表里只留日期、目标和预测变量。
  \item \textbf{缺失值。}目标缺失的行删掉。数值型预测变量的缺失值用该列中位数补上，
        类别型的用众数。
  \item \textbf{日历变量。}新增四列：\code{date\_unix}（自 1970 年起的秒数，即趋势）、
        \code{day\_julian}（一年中的第几天，即季节）、\code{weekday}（1~=~星期一 到
        7~=~星期日，类别型）和 \code{hour}。模型只用列在预测变量里的那几个；窗口总是把
        四个都列上。
  \item \textbf{训练/测试划分。}\code{set} 列把每一行标成 \code{training} 或
        \code{testing}（第~\ref{sec:split}~节）。
\end{enumerate}

\subsection{划分}\label{sec:split}
模型拟合得好不好，由留出的行来评判。划分有四种：
\begin{description}\nosep
  \item[\code{random}]（默认）随机留出 25\,\% 的行。每个季节、星期几和小时在训练集里
        都有代表。
  \item[\code{ts}] 留出记录最后 25\,\%。这时趋势变量必须外推到训练中没见过的范围，
        树模型做不到；用它看预测能力，别用它为归一化训练模型。
  \item[\code{month\_ts}] 在每年的每个日历月里留出连续的一段 25\,\%，它在月中的位置
        随机（带种子）。
  \item[\code{season\_ts}] 同上，只是按气象季节分（DJF、MAM、JJA、SON；12 月归入
        下一年的冬季）。
\end{description}

\caution{随机划分下，相邻的小时落在划分两侧，自相关的记录因此显得比实际好预测：测试集的
$R^2$ 只是一个上界。对若干小时内保持不变的预测变量，这一点最要紧，例如每六小时释放一次的
后向轨迹特征——树模型会拿它当时间戳使。在自带的 MY1 \PMtwo{} 案例上，只用轨迹的模型随机
划分得分约 0.84，分块划分只有 0.46。这类模型请用 \code{month\_ts} 来评判。}

\section{训练模型}\label{sec:train}
\nm{} 默认训练梯度提升树，把预测变量映射到目标。后端有两种：
\begin{description}\nosep
  \item[FLAML] 在墙钟\emph{预算}内，对所勾选的估计器做 AutoML 搜索（默认只搜
        LightGBM；也可加入 XGBoost、限深 XGBoost、随机森林、极端随机树和 scikit-learn 的
        直方图梯度提升），目标是最大化 $R^2$。库的默认预算为 90~s，窗口为 60~s。
  \item[LightGBM] LightGBM 加固定\emph{次数}的随机超参数搜索（窗口里是 20 次）。
\end{description}
预算加长很少有害，对长记录往往有益。FLAML 按墙钟计时，同一脚本、同一数据，在更忙的机器上
可能停在不同的候选模型上；要以次数而不是时间设界，请在 \code{model\_config} 里设
\code{max\_iter}，并把 \code{time\_budget=-1}。

\subsection{模型统计量}\label{sec:modstats}
\code{modStats} 分别对训练行、测试行和全部行做预测，逐组报告：$n$；FAC2（预测落在观测
两倍以内的比例）；MB 和 MGE（平均偏差、平均总误差）；NMB 和 NMGE（同样两项，按观测值之和
归一化）；RMSE；$r$ 及其显著性；COE（效率系数）；IOA（一致性指数）；$R^2$。评判模型，
看留出（$\text{set} = \text{testing}$）那一行。

\section{气象归一化}\label{sec:normalise}
方程~\eqref{eq:norm} 用蒙特卡洛方法算。$N$ 次抽样（\code{n\_samples}，默认 300）中，
每一次：
\begin{enumerate}\nosep
  \item 每一行的\emph{重采样变量}整组换成从\emph{重采样池}（默认就是记录本身）里有放回
        随机抽的一行的对应值。抽的是整行，重采样变量之间的联合结构因此保住了：一次抽样
        绝不会把冬季的温度配上夏季的边界层。
  \item 其余每个预测变量——默认是四个日历变量——保持观测值。
  \item 模型对每一行做预测。
\end{enumerate}
归一化序列就是每个时刻 $N$ 次预测的平均；设了 \code{return\_quantiles}，还会返回 $N$ 次
预测的指定分位数，列名为 \code{q050}、\code{q950} 等（分位数 $\times 1000$）。这些分位数
刻画的是预测在不同天气下的离散程度，不是模型的不确定性。

\paragraph{重采样哪些变量。}重采样气象、保留日历，是通常意义上的去气象化。若连某个日历变量
也重采样，就去掉了它所承载的东西：重采样 \code{day\_julian} 去掉季节循环，重采样
\code{date\_unix} 去掉趋势。非气象的预测变量——比如交通流量——不勾就保持观测值。

\paragraph{抽样多少次。}每个时刻的蒙特卡洛误差按 $1/\sqrt{N}$ 缩小。对平滑的日均或月均，
300 次绰绰有余；单个小时的噪声大些。\code{normalise\_auto} 替你定 $N$：每次加 100 次，
直到连续三次检查都有 95\,\% 的时刻蒙特卡洛相对标准误低于 3\,\%，或到 5000 次为止（到顶时
会警告）。

\paragraph{参考池。}\code{resample\_df} 从另一张表里抽天气——例如用 2019 年的天气，问 2020 年
在 2019 年那样的气象下会是什么样子。\code{conditional\_on} 把池限制在满足条件的行上——一个
值、一组值，或该列上的一个函数，例如 \code{\{"ws": lambda v: v < 2\}} 只从静风时段抽。
\code{resample\_pools} 让不同变量各抽各的池：\code{\{"transport": clean\_hours[traj\_cols]\}}
从清洁气团时段的轨迹特征里抽，本地天气仍旧从整条记录抽。每个池有自己的随机数流，增删或改名
一个池不会动到其他池。

\paragraph{可重复性。}抽样带种子（默认 \code{seed=7654321}），所以模型、数据、设置相同，
序列就逐位相同。

\section{排放分解}\label{sec:emission}
排放分解通过\emph{嵌套}归一化，把归一化序列拆到日历变量所承载的各个时间尺度上：同一粒
种子跑五次归一化，每次多固定一个日历变量，顺序固定如下：
\begin{center}\small
\begin{tabular}{@{}llll@{}}
\toprule
运行 & 保持观测值 & 重采样 & 得到 \\
\midrule
$b$         & 无                                       & 全部       & 基准 \\
$u$         & \code{date\_unix}                        & 其余       & 趋势 $= u - b$ \\
$j$         & \ldots{} + \code{day\_julian}            & 其余       & 季节 $= j - u$ \\
$w$         & \ldots{} + \code{weekday}                & 其余       & 星期 $= w - j$ \\
$h$         & \ldots{} + \code{hour}                   & 气象       & 日变化 $= h - w$ \\
\bottomrule
\end{tabular}
\end{center}
最后一次运行 $h$ 就是普通的归一化序列，报告为 \code{emi\_total}。基准运行拆成均值
\code{emi\_base} 和相对该均值的偏差 \code{emi\_noise}；于是
\begin{equation}
  \text{emi\_total} = \text{emi\_base} + \text{date\_unix} + \text{day\_julian}
     + \text{weekday} + \text{hour} + \text{emi\_noise}.
\end{equation}
不在模型预测变量里的日历变量没有对应的列。

\caution{顺序决定了每个分量装得下什么。算趋势时，一年中的第几天、星期几和小时仍在被平均，
所以趋势装不下每年同一日期重复出现的信号——只能是真正的漂移。算季节分量时，每一行自己的
趋势值已经固定，所以它不是气候态：它仍绑定在所在的那一年上，名字叫“季节”，也可能记下
一次性的东西，比如某年的假期低谷或一次封控。要找重复出现的日历效应，看
\code{day\_julian}，不要看 \code{date\_unix}。}

\section{气象分解}\label{sec:metdecomp}
气象分解把预测中由天气解释的那部分，摊到各个气象变量或它们的命名分组上。日历始终保留观测
值。记 $v(S)$ 为把 $S$ 中变量保留观测值、其余重采样时的归一化序列；$v(\emptyset)$ 即
\code{emi\_total}，$v(\text{all})$ 是不做任何重采样时模型的预测。差值
$v(\text{all}) - v(\emptyset)$ 有两种摊法：
\begin{description}
  \item[顺序归因] 一次固定一个变量，报告每一步的变化：
        $\phi_i = v(\{1,\dots,i\}) - v(\{1,\dots,i-1\})$。$k$ 个变量要 $k+1$ 次归一化，
        但每项贡献都以先前固定的变量为条件，摊法因此取决于顺序——显式给的
        \code{variable\_order}，没给就用模型的特征重要性，而模型重拟合时它可能变。
  \item[Shapley 归因] 对每个变量，在所有可能的固定顺序下取边际效应的平均：
        \begin{equation}
          \phi_i = \sum_{S \not\ni i} \frac{|S|!\,(k-|S|-1)!}{k!}\,
                   \bigl[v(S\cup\{i\}) - v(S)\bigr].
        \end{equation}
        没有哪个顺序占便宜。精确值要全部 $2^k$ 个联盟，最多容得下 $k=10$；再多就把
        变量分组，或传入 \code{n\_permutations}，用随机顺序（成对的对偶抽样）来估计。
\end{description}
无论哪种，各项贡献之和严格等于 $v(\text{all}) - v(\emptyset)$；每次归一化用同一粒种子，
所以它们是配对的差（公共随机数），不是两个含噪估计相减。结果里还有
\code{met\_total} = 观测 $-$ \code{emi\_total}、它的均值 \code{met\_base}，以及
\code{met\_noise} = \code{met\_total} $-$ \code{met\_base} $-\sum_i\phi_i$，即观测中模型
没有再现的部分。

\paragraph{分组。}\code{groups=\{"local": met\_cols, "transport": traj\_cols\}} 给每个分组
一列；分组作为整体固定、整体重采样，所以它的列含着组内成员之间的交互作用。每个气象预测
变量必须恰好落在一个分组里。有分组时默认 Shapley，没有分组时默认顺序归因。

\section{滚动归一化}\label{sec:rolling}
\code{rolling} 逐窗口归一化整条记录：每 \code{rolling\_every} 天（7）起一个长
\code{window\_days} 天（默认 14）的窗口，每个窗口的天气\emph{只从窗口内}抽。所以短窗口
只去掉在窗口内变化的那部分天气——两周的窗口就是天气尺度的变化——而保留全记录归一化会去掉
的较慢气象变化。结果每个窗口一列，从 1 编号：\code{window\_01}、\code{window\_02}、……
（超过 99 个窗口时补成三位，列名便按窗口顺序排）；\code{window\_columns(rolled)} 列出这些
列。窗口里显示它们的均值，以及重叠窗口之间的离散度。

\section{多尺度分解}\label{sec:multiscale}
记 $Y_W$ 为窗口宽 $W$ 下滚动归一化序列的均值，$Y_\infty$ 为普通归一化序列（池是整条
记录）。$W$ 越大，窗口的池越接近整条记录，$Y_W \to Y_\infty$。不同宽度之间的差，把每个
时间尺度频带特有的天气影响分离出来，道理和小波细节系数一样：
\begin{align}
  D_\text{fast} &= Y_\text{fast} - Y_\text{meso} && \text{天气尺度，月以下} \\
  D_\text{meso} &= Y_\text{meso} - Y_\text{slow} && \text{季节内} \\
  D_\text{slow} &= Y_\text{slow} - Y_\infty      && \text{连慢窗口的池也留下的部分}
\end{align}
宽度默认 14、90 和 365~天。每个尺度的窗口以宽度的四分之一为步长；记录上不足两个窗口的
尺度会被跳过，并给出说明。

\section{部分依赖}\label{sec:pdptheory}
部分依赖曲线展示模型的预测如何随一个预测变量变，其他变量在数据上取平均：在该变量第 1 到
第 99 百分位之间取 \code{grid\_points} 个值，对每个值把每个训练行的该变量都设成它，再对
预测取平均。\code{pdp\_mean} 就是这个平均，\code{pdp\_std} 是预测绕它的离散度。曲线描述
的是模型，不是大气；相关的预测变量（比如温度和边界层高度）一次只动一个，会走到记录里可能
从没出现过的组合。

\section{用 Chronos-2 零样本去气象化}\label{sec:chronos}
Chronos-2 是预训练的时间序列基础模型。\nm{} 直接拿来用，不经训练：气象走模型的协变量
通道，天气则照方程~\eqref{eq:norm} 原样通过重采样边缘化。历史短到不足以训练、刚投运的
站点，适合用它；或者把它当作一个从没见过目标自身记录的模型，给出第二个估计。四个特性决定
它合不合用：
\begin{itemize}\nosep
  \item \textbf{前 \code{context\_length} 行不算结果。}模型要这么长的历史（默认 2048 行）
        才能外推，所以这些行用观测值填，两条曲线在那里按构造重合。给它几年的数据，别只给
        几个月。
  \item \textbf{每次蒙特卡洛抽样都是一次完整的前向计算}，所以默认只抽 8 次，不是 300 次。
        模型按优先级在 CUDA、Apple Silicon（\code{mps}）或 CPU 上跑；CPU 很慢。
  \item \textbf{位置即时间。}Chronos-2 数的是行，记录必须落在规则网格上；
        \code{to\_regular\_index} 会重建网格，空缺留作缺失。
  \item \textbf{协变量敏感性是质量门槛。}没有拟合任何东西，就没有留出集拟合，也没有特征
        重要性。取而代之的是打乱未来的天气，量预测中位数偏移多少。若低于预测值的约
        1\,\%，说明模型在读目标自身的历史、不理天气，它给的“归一化”序列没有意义。
\end{itemize}
气象分解有零样本路径。排放分解、滚动归一化和部分依赖没有：它们定义在已拟合模型对日历变量
的响应上，而 Chronos-2 是从目标自身的历史里读出这些的。

\section{合成控制}\label{sec:scmtheory}
合成控制用一份在\emph{干预日期}前后都有观测的单元面板，估计一项干预对一个\emph{处理}单元
（一个站点、一座城市）的效应。干预没碰过的单元——\emph{供体池}——组合成处理单元的合成
版本，在干预日期之前跟住它；干预日期之后，合成序列就是反事实，而
\begin{equation}
  \text{效应}_t = \text{观测}_t - \text{合成}_t .
\end{equation}
干预后时段的平均效应报告为 ATT（处理组平均处理效应）。就空气质量结果而言，去气象化的序列
是自然的输入：这样天气就不属于供体必须再现的东西。

\subsection{估计器}
\begin{description}\nosep
  \item[\code{scm}] 增广 SCM（默认）。在每个日期对供体做岭回归，用各单元干预前的轨迹
        预测结果变量；合成值是处理单元的岭回归预测加上供体岭回归残差的加权和，权重之和为
        一（除非放开，否则非负）。
  \item[\code{abadie}] 经典 SCM：非负、和为一的供体权重，拟合处理单元干预前的路径。
        好解释，但处理单元落在供体范围之外时拟合不了。
  \item[\code{did}] 双重差分：假设平行趋势的双向固定效应基准。
  \item[\code{mcnnm}] 核范数最小化的矩阵补全（Athey 等，2021）：把干预后处理单元的格子
        当作缺失，用面板的低秩拟合补上。
  \item[\code{robust}] 稳健 SCM（Amjad、Shah 和 Shen，2018）：拟合权重前，先用奇异值
        阈值给供体矩阵去噪。
  \item[\code{mlscm}] ML 增广 SCM（\emph{实验性}）：一个 AutoML 模型在干预前时段把供体的
        结果映射到处理单元的结果。它没有其他方法的因果保证，并会就此发出警告。
  \item[\code{bayesian}] 贝叶斯单纯形 SCM（PyMC）：权重取 Dirichlet 先验，后验可信区间
        直接给出。
\end{description}

\subsection{诊断}\label{sec:scmdiag}
\code{scm\_diagnostics} 报告干预前的拟合（\code{pre\_rmse}、\code{pre\_mae}、
\code{pre\_mape}、\code{pre\_r2}）、干预后的效应（\code{att}、\code{att\_cum}、
\code{post\_rmse}），以及权重的集中程度：Herfindahl 指数 $\text{HHI} = \sum_j w_j^2$、
有效供体数 $1/\text{HHI}$、最大权重的供体及其份额。

\caution{干预前拟合得好，只有当估计器不能把任意序列都拟合得一样好时，才算证据。凡是靠供体
加权去匹配干预前的估计器，供体数相对干预前日期数一多，就会过拟合，增广 SCM（\code{scm}）
最厉害：它的岭回归特征就是各单元干预前的结果，所以在某个干预前日期，处理单元自己当天的值
就是用来预测它的特征之一；供体数和日期数一样多时，它几乎能再现任何干预前时段。第\ref{ch:scm}
章的 ULEZ 面板上有 27 个干预前月份：把处理序列换成随机噪声，5 个供体拟出 $R^2 = -0.05$，
20 个得 0.69，53 个得 0.997；经典 SCM 拿 40 个随机游走供体拟合 12 个月，打乱的序列也能得
0.60。因此 SCM 工作室会把处理单元干预前的值按时间打乱，用同一设计再拟一次，只有打乱序列
达不到 0.5 时，才解读干预前的 $R^2$（第~\ref{sec:scmfit}~节）。\code{scm} 估计器的供体权重
只平衡岭回归残差，所以它们集中不集中，并不说明合成单元由哪些供体构成。}

\subsection{推断}\label{sec:scminference}
\begin{description}
  \item[空间安慰剂] 依次把每个供体当作处理单元，用\emph{其他}供体为它构造合成控制（真正的
        处理单元从不做安慰剂的供体）。以干预后平均效应为统计量，置换 $p$ 值为
        \begin{equation}
          p = \frac{1 + \#\{\,j : |\bar{e}_j| \ge |\bar{e}_\text{treated}|\,\}}{1 + J},
        \end{equation}
        其中 $J$ 是安慰剂单元数；最小只能到 $1/(J+1)$。安慰剂效应还给处理效应提供逐点
        区间。
  \item[时间安慰剂] 把干预日期挪到更早、什么都没发生的日期（每 \code{placebo\_every} 步
        一个，前面至少留 \code{min\_pre\_period} 步）。真实效应按同样的方式与这些伪干预
        日期的效应排序。每个伪干预后时段都必须在真实干预日期之前结束。
  \item[刀切法区间] 依次去掉每个供体重新拟合，效应的离散度给出逐点区间。靠单个供体撑起来
        的结果会在这里现形。
  \item[自助法区间] 在供体的随机子集上（默认 80\,\%）重复拟合 $B$ 次，可选对干预前时段
        做分块重采样。
  \item[共形区间] \code{conformal\_effect\_interval} 给出干预后平均效应的有限样本区间
        （Chernozhukov、W\"uthrich 和 Zhu，2021）；\code{rmspe\_ratio\_test} 给出 Abadie
        的干预后/干预前 RMSPE 比值检验。
\end{description}

\subsection{交错采用}\label{sec:scmstaggered}
干预分批落地时——伦敦 ULEZ 2021 年扩区、限行政策在各城市错峰实施——每个采用单元有自己的
干预日期。\code{run\_scm\_staggered} 为每个处理单元按各自的干预前/后时段单独拟合，并把
\emph{其余处理单元}一律排除在每次拟合的供体池之外，一个单元的干预便漏不进另一个单元的
合成控制。每个单元的效应是其效应路径在干预后的平均，
\begin{equation}
  \bar{e}_j = \frac{1}{|T_j|} \sum_{t \ge c_j} \bigl(\text{观测}_t - \text{合成}_t\bigr),
\end{equation}
其中 $c_j$ 是该单元的干预日期、$T_j$ 是其干预后日期。结果每个单元一行（$c_j$、干预前/后
长度、干预前拟合 $R^2$、效应，以及 $p_j$——该单元在自身干预日期处的空间安慰剂检验），再
加一行各单元效应的均值/中位。

\emph{均值}行的 $p$ 值是跨单元置换检验：每个供体依次扮演每个处理单元的角色——全套干预
日期——其伪效应按与真实单元相同的方式平均，
\begin{equation}
  p = \frac{1 + \#\{\,u : |\bar{e}_u| \ge |\bar{\bar{e}}|\,\}}{1 + J}，
\end{equation}
其中 $J$ 为供体数、$\bar{\bar{e}}$ 为各处理单元效应的均值。真实处理单元不进参考分布：
它们的数据已含干预——交错设计里，给某单元一个比它真实日期更晚的干预日期，它的真实效应就
落在新的「干预前时段」里；给更早的日期，「干预后时段」又部分无效应。被置换的处理身份被
真实效应污染，且处理单元很少时检验退化——两个单元只有一个非恒等置换，且被污染。只有供体
的参考是干净的零假设：从未受干预的供体给出零假设下的抽样。每个安慰剂拟合需要自身可用的
供体池，默认 \code{scm} 后端至少要 3 个供体；更小的池 $p$ 值为 NaN，但效应仍读出。

% ============================ 4. GETTING STARTED ==============================
\chapter{入门}\label{ch:start}

\section{打开表格之前的窗口}
\nm{} 启动时停在\tab{数据}标签页，什么数据都还没载入（图~\ref{fig:welcome}）。参数面板收着，
助手（第\ref{ch:assistant}章）在结果旁边打开。操作栏上的\emph{展开面板}或 Ctrl+B 展开参数
面板；\tab{视图 $\rightarrow$ 助手}把助手关掉。\emph{从时间序列开始}下面有两个按钮——
\emph{打开数据…}和\emph{询问助手}——再往下是几张入口卡片：示例；最近打开的五张表（各附
文件夹、大小和最后修改日期，之后被挪过地方的会变灰）；还有两处不用自备文件就能拿数据的地方
——数据工作室和传输工作室。没内容可列的卡片会省略。把表格文件（第~\ref{sec:input}~节）拖到
窗口任意位置就能打开。

\screenshot{ui_welcome.png}{启动时的窗口：\tab{数据}标签页上的欢迎页，标签页上方是操作栏
（第一个按钮为\emph{展开面板}），右侧是助手，底部是日志。参数面板收着，需要时再展开。}{fig:welcome}

\section{示例}\label{sec:examples}
示例在\tab{文件 $\rightarrow$ 载入示例}和欢迎页的\emph{示例}卡片上；打开任一示例时，目标和
预测变量都按它来源的设定勾好。
\begin{description}
  \item[伦敦 Marylebone Road \NOtwo{}（2020）] \nm{} 论文的去气象化案例：一个路边站点的
        逐时 \NOtwo{} 和 ERA5 气象，2020 年 1 至 8 月，其间有英国的 COVID-19 封控（即下表
        的 \code{load\_my1()}）。本手册主窗口的图都由它生成。
  \item[伦敦 Marylebone Road \PMtwo{}（2020）] 论文中考虑传输的案例：同一站点的逐时
        \PMtwo{}，配本站的风和温度以及 ERA5 的其他字段（\code{load\_my1\_pm25()}）。
  \item[合成站点，一年逐时] 即时生成的一年逐时合成数据：\PMtwo{}、PM\textsubscript{10}、\NOtwo{}、O\textsubscript{3}、
        SO\textsubscript{2} 和 CO，另有 ERA5 风格的气象（\code{u10}、\code{v10}、
        \code{t2m}、\code{d2m}、\code{blh}、\code{sp}、\code{ssrd}、\code{tcc}、\code{tp}、
        \code{rh2m}、\code{ws}、\code{wd}）。污染物里混了季节循环、早晚高峰循环和对风速的
        依赖，每一步都有东西可看。目标预设为 \PMtwo{}，气象已勾好。
  \item[合成控制面板] 打开 SCM 工作室并载入它自己的示例；SCM 工作室的\tab{文件
        $\rightarrow$ 载入示例}下有两个（第\ref{ch:scm}章）：论文里的伦敦 ULEZ 面板，以及
        一个 31 单元的合成周面板，其中单元“2+26 cities”的 SO\textsubscript{2} 自
        2015-10-23 起被削减约 35\,\%，其余单元都是干净的供体。
\end{description}
论文里的真实记录随包提供，脚本中这样载入：
\begin{center}\small
\begin{tabular}{@{}lp{7.0cm}p{3.2cm}@{}}
\toprule
载入函数 & 内容 & 用于 \\
\midrule
\code{load\_my1()}      & 伦敦 Marylebone Road 逐时 \NOtwo{} + ERA5 气象，2020 年 1--8 月 & 去气象化 \\
\code{load\_scm()}      & 104 个英国站点的月均去气象化 \NOtwo{}，2016--2021 & 合成控制（ULEZ） \\
\code{load\_my1\_pm25()} & MY1 逐时 \PMtwo{} + 气象，2020 年 1--8 月 & 考虑传输的案例 \\
\code{load\_traj\_my1()} & 到达 MY1 的 6 小时一次的 HYSPLIT 轨迹特征 & 考虑传输的案例 \\
\code{example\_traj\_dir()} & 两天的原始 HYSPLIT \code{tdump} 文件 & 传输工作室 \\
\bottomrule
\end{tabular}
\end{center}
数据来源：英国 AURN（Defra，Open Government Licence v3.0）；ERA5（Copernicus 气候变化服务）；
HYSPLIT（NOAA ARL）。

\section{第一次分析}
\begin{enumerate}\nosep
  \item \tab{载入数据。}\emph{打开数据…}，选一张含时间、污染物和气象的表
        （第\ref{ch:data}章）；设好站点用的\emph{时区}（第~\ref{sec:clock}~节）。用\emph{展开面板}（Ctrl+B）展开参数面板，或直接问助手。
  \item \tab{选择列。}检查目标，勾选预测变量（第~\ref{sec:columns}~节）。
  \item \tab{训练}，读\tab{模型}标签页：留出集上的拟合够不够用来归一化？
        （第\ref{ch:train}章）
  \item \tab{归一化}，把归一化序列与观测序列比一比（第\ref{ch:normalise}章）。
  \item 按问题需要做\tab{分解}、\tab{滚动}和\tab{部分依赖}
        （第\ref{ch:decompose}--\ref{ch:pdp}章）；\tab{趋势}检验回答分解刚画出的
        序列是在升、在降还是平的（第~\ref{sec:trend}~节），\tab{不确定性}标签页回答
        区间的离散来自哪里（第~\ref{sec:uncertainty}~节）。
  \item \tab{保存结果：}导出表格、HTML 报告，把会话存成脚本（第\ref{ch:export}章）。
\end{enumerate}

% ============================ 5. THE INTERFACE ================================
\chapter{界面}\label{ch:ui}

主窗口左边是\emph{参数面板}，右边是\emph{结果标签页}。参数面板按工作流程的每一步分成一页，
从面板顶部的标签栏选——\tab{数据}、\tab{模型}、\tab{归一化}、\tab{分解}、\tab{滚动}和
\tab{PDP}——每一步都在决定它的设置旁边带一个自己的\emph{运行}按钮。结果落在同名步骤的
标签页里，该步骤一完成，这个标签页就切到前面。参数面板启动时收着，把宽度让给结果和助手；
操作栏上的\emph{折叠 / 展开面板}（Ctrl+B）收放它。面板收起时，窗口能缩到约 740 像素宽；
再窄，操作栏就折成两行。

\section{参数面板}
各页及其中的内容（每页第一个框不加标题，由标签名说明）：
\begin{description}\nosep
  \item[数据] \emph{打开数据…}、文件名及其行数列数、时间用的\emph{时区}，以及时间是
        怎么读出来的（第~\ref{sec:input}~节）；再往下是\emph{列}：目标和预测变量
        （第~\ref{sec:columns}~节）。
  \item[模型] 后端、划分、训练集比例、时间预算（或搜索次数）、估计器、随机种子；换到
        Chronos-2 则是设备和微调（第\ref{ch:train}章）。
  \item[归一化] 重采样哪些变量、样本数、CPU 核心数、分位数带（第\ref{ch:normalise}章）。
  \item[分解] 排放或气象；气象分解另有归因方式和分组（第\ref{ch:decompose}章）。往下
        依次是\emph{多尺度分解}：三种窗口宽度（第~\ref{sec:multiscaleui}~节）；\emph{趋势}：
        要检验的序列（分解的分量、它分离出的长期趋势，或观测目标）、季节周期、段数、
        自助法重采样次数；\emph{事件与异常}：要打分的序列、方法、阈值和最短标记长度
        （第~\ref{sec:eventsui}~节）；\emph{不确定度分解}：成员数、各成员共用的气象抽样数、
        成员后端（第~\ref{sec:uncertainty}~节）。
  \item[滚动] 窗口、步长和样本数（第~\ref{sec:rollingui}~节）。
  \item[PDP] 变量和网格点数（第\ref{ch:pdp}章）。
\end{description}
参数面板的宽度取最宽的那一页。面板的配置从\tab{文件}菜单保存和载入
（第~\ref{sec:config}~节）。

\section{结果标签页}
\begin{center}\small
\begin{tabular}{@{}lll@{}}
\toprule
标签页 & 显示 & 由谁填充 \\
\midrule
数据     & 概览、周期、风玫瑰和污染玫瑰、统计；排除的行 & 载入 \\
模型     & 评估结论、统计表、散点对比图、特征重要性、前向滚动交叉验证 & 训练模型；前向滚动交叉验证 \\
归一化   & 观测序列和归一化序列、分位数带及其周期 & 运行归一化 \\
分解     & 排放分量或气象分量，以及周期 & 运行分解 \\
滚动     & 各滚动窗口的均值及其离散度 & 运行滚动归一化 \\
多尺度   & 尺度空间级联和各频带的能量 & 多尺度 \\
趋势     & 被检验的序列及其 Theil--Sen 直线，或其各段折线与断点、评估结论和 p 值 & 趋势分析 \\
事件     & 异常分数条，被标记的窗口以阴影显示，及其事件表 & 查找事件 \\
PDP      & 每个变量一幅部分依赖图 & 计算 PDP \\
不确定性 & 逐时刻的离散按来源堆叠（气象、模型形式、二者的交互）及中位数份额 & 不确定度分解 \\
\bottomrule
\end{tabular}
\end{center}
需要做判断的标签页开头是一条\emph{评估横幅}：绿色（好）、琥珀色（注意）或红色（不可用），
后面几行文字说明运行了什么，有表格时再附一张简表。每幅图顶部都有一条工具栏，图标小巧、与
窗口配色一致——主页、后退、前进、平移、缩放、子图间距、坐标轴和曲线选项、保存——然后是
\emph{打印用图…}（第~\ref{sec:print}~节），最右端显示指针所在处的数值。不管窗口是什么
外观，图都以白底深色保存。每幅图都按期刊要求绘制：7 到 8.5~pt 的文字，细坐标轴，刻度朝外，
图例不加边框，色盲读者也分得清的配色，化学式带下标（\NOtwo{}、\PMtwo{}），日期轴每一级
只写一次——先写年份，再写月份——共用一条时间轴的堆叠子图只在最下面一幅标日期。

\paragraph{时间分辨率。}\tab{数据}、\tab{归一化}、\tab{分解}、\tab{滚动}和\tab{多尺度}
标签页上的时间序列图上方，有一个\emph{时间分辨率}选择器：\emph{原始}、\emph{逐时}、
\emph{日均}、\emph{周均}、\emph{月均}或\emph{年均}。选更粗的步长，图就按每个时段的均值
重画，长的逐时记录看起来省力；比记录本身还细的步长呈灰色。比逐日更细、点数超过 2,000 的
记录默认按日均显示——几个月的逐时数据原样画出来只是一条实心的带——在某个标签页上做的选择
会沿用到它的下一个结果。变的只是显示：结果、结果表及其导出仍保持记录本身的分辨率。记录里
有空缺（间隔超过通常步长的五倍）时，曲线在空缺处断开，不用直线连过去。图里给出的 $\sigma$
总是按记录本身算，与显示的步长无关。

同一个模型上排放分解和气象分解都跑过时，\tab{分解}标签页两个都留着：\emph{时间分辨率}旁边
的\emph{方法}选择器在两者之间切换，保存的会话也会把两者都恢复。

\section{操作栏}
标签页上方依次是：\emph{折叠 / 展开面板}；并排的几个工作室——\emph{获取数据}（数据工作室，
第~\ref{sec:datastudio}~节）、\emph{传输工作室}（第~\ref{sec:transport}~节）和\emph{SCM
工作室}（第\ref{ch:scm}章）；然后是对结果的操作——\emph{导出结果 CSV}（当前标签页的表）和
\emph{HTML 报告}（当前标签页的结果连同其溯源信息，做成报告）。

\section{运行按钮}
\begin{center}\small
\begin{tabular}{@{}lp{5.0cm}l@{}}
\toprule
按钮 & 需要 & 运行 \\
\midrule
训练模型 / 载入 Chronos-2 & 一张表，至少勾选一个预测变量 & \code{build\_model} \\
微调 Chronos-2        & 已载入 Chronos-2         & \code{finetune} \\
运行归一化            & 一个模型，至少一个重采样变量 & \code{normalise} \\
运行分解              & 一个模型                 & \code{decompose} \\
运行滚动归一化        & 一个 AutoML 模型         & \code{rolling} \\
多尺度分解            & 一个 AutoML 模型，一次归一化 & \code{rolling} $\times 3$ \\
趋势分析              & 一个 AutoML 模型（有分解时检验分解的序列） & \code{trend}（1 段）/\code{detect\_trend\_breaks}（更多） \\
查找事件              & 一个模型（给分量打分时须有分解） & \code{anomaly\_scores}、\code{detect\_events} \\
计算 PDP              & 一个 AutoML 模型，至少勾选一个变量 & \code{pdp} \\
不确定度分解          & 一个 AutoML 模型 & \code{decompose\_uncertainty} \\
\bottomrule
\end{tabular}
\end{center}
按钮只在能执行时可用；不能执行时，工具提示会说明缺什么。一次只跑一个任务，在后台线程上跑，
窗口所以一直能响应；运行期间各运行按钮禁用，\emph{取消}可用，状态栏显示任务名和已用时间。
任务能报进度时——滚动窗口、分解的各步、LightGBM 的各次试验、HYSPLIT 的各个受体时刻、FLAML
的时间预算——进度条跟着走，状态栏给出完成比例，依据够了还会给剩余时间；报不出进度的任务，
进度条是忙碌条。任务结束时，这一步写出的提示会说明用了多久。取消会\emph{放弃}该任务：窗口
立刻空出来，但计算停不下来，它会在后台跑完，结果丢掉。每一步的调用成功时都进命令日志
（第~\ref{sec:commandlog}~节）。

前提还不具备就要跑某一步——模型还没训练就归一化，归一化之前就跑多尺度——窗口不弹对话框，
而是在结果上方显示一条\emph{通知}：说明缺什么，并把能解决它的那一步放在按钮上（\emph{训练
模型}、\emph{运行归一化}）；\(\times\) 可以关掉它，那一步做完它也会自己消失。状态栏右侧每
一步有一盏灯——数据、模型、归一化、分解、多尺度、趋势、事件、不确定度、滚动、PDP——那一步
有了结果就变绿。

\section{菜单和快捷键}
\begin{center}\small
\begin{tabular}{@{}lll@{}}
\toprule
菜单 & 命令 & 快捷键 \\
\midrule
文件     & 打开数据…                            & Ctrl+O \\
         & 载入示例（各示例、SCM 的示例）       & \\
         & 获取数据（空气质量网络 + 气象）…     & Ctrl+D \\
         & 最近打开（八张表）                   & \\
         & 保存结果… / 载入结果…                & Ctrl+S \\
         & 重新打开上次会话                     & \\
         & 载入会话脚本…                        & \\
         & 方法段落…                            & \\
         & 导出当前结果…                        & Ctrl+E \\
         & 导出 HTML 报告：English、中文         & \\
         & 将会话保存为脚本…                    & \\
         & 保存配置… / 载入配置…                & \\
         & 运行历史…                            & \\
视图     & 命令面板…                            & Ctrl+K \\
         & 折叠 / 展开面板                      & Ctrl+B \\
         & 日志…、命令日志…                     & \\
         & 助手                                 & Ctrl+Shift+A \\
         & 外观：跟随系统、浅色、深色           & \\
         & 语言：跟随系统、English、中文        & \\
分析     & 合成控制（SCM）工作室…               & Ctrl+Shift+S \\
         & 风向玫瑰图…                          & \\
帮助     & 用户手册                             & F1 \\
         & 键盘快捷键…                          & \\
         & 文档、关于 Normet                    & \\
\bottomrule
\end{tabular}
\end{center}
在 macOS 上，Ctrl 就是 Command 键。\emph{命令面板}（Ctrl+K）列出每一条菜单命令和每一个
结果标签页；敲几个字母再回车就行。\tab{帮助 $\rightarrow$ 键盘快捷键}从菜单里读出来，逐条
列出每个快捷键；\tab{帮助 $\rightarrow$ 用户手册}（F1）按界面语言打开本手册。

\section{日志和助手}
\tab{视图 $\rightarrow$ 日志…}在单独的窗口里打开 \nm{} 的日志，可以一直开在主窗口旁边：
每一步做了什么、多少行、FLAML 选中了哪个估计器及其超参数、一次滚动跑出了多少窗口。
\tab{视图 $\rightarrow$ 命令日志…}打开能重复本次会话的调用（第~\ref{sec:commandlog}~节）。
任务失败时，完整回溯写进日志，最后一行显示在对话框里。\emph{助手}占住窗口右侧的整个高度
（第\ref{ch:assistant}章）。

\section{外观和语言}\label{sec:appearance}
窗口和其中的图跟随系统的浅色或深色外观；\tab{视图 $\rightarrow$ 外观}可以固定用其中一种，
保留到下次会话。深色窗口上的图只为显示而重新着色——底色、文字、刻度和网格；承载数据的颜色
不动——保存时用浅色。

窗口可以显示英文或中文：默认跟随系统语言，也可以在\tab{视图 $\rightarrow$ 语言}里选，选择
保留到下次启动。切换立即生效——菜单、按钮、标签页、工具提示、评估结论，以及图的标题、坐标轴
标签和月份名，都用所选语言重新写出，图按显示的样子保存。数据自身的文字（列名、站点代码）、
日志、命令日志里的代码和助手的对话保持原样。系统自己画的部分则不论窗口用哪种语言都保持系统
语言：打开和保存面板连同它们的标题，以及 macOS 自己加进菜单的项（\emph{自动填充}、\emph{开始
听写}、\emph{表情与符号}）。HTML 报告（第~\ref{sec:commandlog}~节）用导出时选的语言写——
\tab{文件 $\rightarrow$ 导出 HTML 报告}或\emph{HTML 报告}按钮列出各语言，每种都用它自己的
文字命名——包括标题、标签和图；所选语言留给下一份报告，没选过就用系统语言。

窗口的大小、位置、面板页和结果标签页，按上次关闭时的状态保留。结果表可点表头排序，右键
复制为制表符分隔的文本。任务在窗口不在前台时完成，会用提示音和 Dock 弹跳把分析者唤回来
（系统允许时还有 macOS 通知）；清空运行历史前会先问一声。步骤灯和助手对话里已完成的工具
调用，除点击外也能用键盘打开。

% ============================ 6. COLUMNS AND TRAINING =========================
\chapter{训练模型}\label{ch:train}

\section{表格及其列}\label{sec:columns}
主窗口打开的表必须有一列 \code{date}（不区分大小写），取值能解析成日期；其余每个数值列都
可以当目标或预测变量。目标预设为第一个名字像污染物的列（\PMtwo{}、PM10、\NOtwo{}、O3、SO2、
CO、NOx），预测变量预设为被认作气象的列：
\begin{center}\small
\code{u10 v10 d2m t2m blh sp ssrd tcc tp rh2m ws wd temp rh pressure}
\end{center}
（匹配不区分大小写）。\emph{全部}把预测变量全勾上——已全勾时再点一次就全取消——\emph{仅
气象}只勾气象变量。目标即便勾上，也从不作为预测变量。四个日历变量自动进每个模型，不在列表
里出现。

\note{换了名字的气象列认不出来，得手动勾——MY1 里就是 \code{atmos\_pres}。\tab{归一化}页
上的\emph{仅气象}也一样（第\ref{ch:normalise}章）。}

\screenshot{ui_data.png}{从\tab{文件 $\rightarrow$ 载入示例}打开 MY1：2020-01-01 至
2020-08-31 共 5,793 个逐时行、11 列，无缺失值。示例设好了论文的目标 \NOtwo{} 及其九个气象
预测变量（其中 \code{atmos\_pres} 的名字窗口不认识）。\tab{概览}页把每一列画成热图（各自
按自身范围缩放），下方是所选的序列（此处为 \NOtwo{} 日均）。}{fig:data}

\tab{数据}标签页有四个子页。\tab{概览}一眼看全整条记录：在图上拖选一段时间，或点选光标下
的点（序列按时间分辨率平均时，选的是该点代表的时段），里面的行列在图下方；按住 Shift（或
Ctrl、Cmd）拖动或点击可以加选，按住 Shift 点已选的点则取消它，按 Esc 清空选择。图下方一行，
左边显示所选内容——行数，其中多少行已排除——右边显示光标所指的时间、变量和数值。列表里可以
像电子表格那样选行——按住 Ctrl 或 Cmd 点击选不相邻的行，按住 Shift 点击或沿行拖动选一段——
再在这些行或其行号上右键排除或恢复（右键落在未选中的行上时，先单独选中它）。在图上右键，
可对全部已选取的行做同样的事，点在选取范围外时先选中那个点；\emph{时间分辨率}一行上的
\emph{排除所选}和\emph{恢复所选}作用于选中的行，一行没选时作用于全部已选取的行。列表留着，
被排除的行显示为灰色，方便接着操作。旁边的\emph{排除}框接受逐行输入的时间段
（\code{2020-03-23 .. 2020-04-05}，或一个日期），按\emph{应用}生效。排除的行在图上打阴影，
并从之后每一步里去掉：训练、归一化、分解等。命令日志把它们写成 \code{left\_out} 和
\code{df\_kept}，会话也保存它们。\tab{周期}画一列的日变化、周变化和季节变化；表里有风速
风向时，\tab{风}画风玫瑰图、污染物的污染玫瑰图和极坐标图；\tab{统计}给出每列的值个数、
缺失比例、均值、标准差、最小值、四分位数和最大值。四个子页都只读保留下来的行。

\section{设置}
\begin{description}\nosep
  \item[后端] \code{flaml}、\code{lightgbm} 或 \code{chronos-2}
        （第~\ref{sec:chronosui}~节）。没装 Chronos-2 的扩展时呈灰色，工具提示说明原因。
  \item[划分方法、训练集比例] 见第~\ref{sec:split}~节；默认 0.75。
  \item[时间预算] FLAML 的搜索时间，默认 60~s。对 LightGBM，这一行变成\emph{搜索次数}，
        默认 20；切换后端时各自的值都留着。
  \item[估计器]（仅 FLAML）默认只搜 LightGBM；加更多以扩大搜索范围。
  \item[随机种子] 默认 7654321：划分和搜索都用它做种子。
\end{description}

\section{阅读模型标签页}
\emph{训练模型}运行 \code{build\_model}——准备数据、划分、搜索——并填充\tab{模型}标签页
（图~\ref{fig:model}）。评估结论看留出集的 $R^2$：0.7 及以上为拟合良好，0.5 到 0.7 为
一般，低于 0.5 为差。表格是训练行、测试行和全部行的 \code{modStats}
（第~\ref{sec:modstats}~节）；散点对比图显示留出行的观测与预测，画成正方形，附 1:1 线，
角上标出行数、预测对观测的回归斜率（小于 1：极值被拉向中间）和平均偏差——行数上千时改用
六边形密度图；条形图是模型的特征重要性（LightGBM 为分裂次数）。表格下方一行给出点预测的
\emph{CRPS 参考值}——留出行上的 $\sqrt{\pi/2}\cdot\mathrm{MAE}$，即点预测在连续分级概率
评分（CRPS）尺度上的得分。带分布的预报算得出自己的 CRPS，只有低于这个参考值，才说明这个
分布值得带。

\screenshot{ui_model.png}{用 FLAML（LightGBM，30~s 预算，随机划分）训练 MY1：在 1,448 个
留出小时上 $R^2 = 0.883$、RMSE $= 8.66$~\ugm{}，而在 4,345 个训练小时上为 $0.976$。
风向风速、边界层高度和趋势承担了最多的分裂。}{fig:model}

提升树在训练行上总比在留出行上拟合得好，这在意料之中，本身不是问题；留出集上拟合差才是
问题。拿一般或差的模型去归一化之前，先试试更多预测变量、更长的预算或更多估计器。归一化
序列会继承模型的误差。

一次随机划分是在给模型化妆：每个测试小时的邻域都在训练行里，自相关的序列靠局部信息就能
预测得很好，但这不代表整体预测得好。\emph{前向滚动交叉验证}正是对这个数字的检验：同一后端
在扩展窗口上重训，给紧接其后的数据段打分，默认做五次（\emph{CV 折数}），每折花一次训练
预算。标签页因此多了一个\emph{视图}选择器，在划分的图与各折之间切换——每折一根 $R^2$ 柱，
划分的取值以横线画在柱上；表里是各折的测试窗口、行数与 $R^2$、RMSE 和平均偏差。划分的
横线要是明显高过柱群，评估结论会直说：请以各折为准。对状态会漂移的序列，随机划分抽未来
和抽过去一样随意，各折的数字尤其可信。

\section{Chronos-2}\label{sec:chronosui}
选择 \code{chronos-2} 时，训练相关的控件藏起来——没有东西要训练，没有预算，没有划分——按钮
变成\emph{载入 Chronos-2}。第一次载入会下载约 500~MB 的权重。\emph{设备}选在哪里跑
（\emph{auto}：先 CUDA，再 Apple Silicon，最后 CPU）。至少要勾一个气象预测变量：没有协变量，
模型做的只是普通预测，不是去气象化。

此后\tab{模型}标签页报告检查点、设备、上下文长度和预测长度，以及协变量敏感性检查
（第~\ref{sec:chronos}~节）：打乱未来天气时预测中位数的平均偏移和最大偏移，以及该偏移占
预测值的百分比。5\,\% 及以上为绿色，1 到 5\,\% 为琥珀色（序列主要由目标自身的历史承载），
低于 1\,\% 为红色（模型在做自回归，它给出的去气象化序列毫无意义）。

表里还有一项分位数回测：检查点对最后一段预测长度做一次 10/50/90 分位预测，用 CRPS 把这个
分布带对观测尾部打分（CDF 由三个分位点线性插值），它的中位数则按点预测打分（还是
$\sqrt{\pi/2}\cdot\mathrm{MAE}$）。表下一行说明这个分布带有没有打败其中位数的点参考值——
即这个分布值不值得带——pinball 损失再按分位逐级拆解。这一行会写清它依据的窗口：一段预测
长度只是快速检查，不是论文里要引用的回测。

\emph{微调 Chronos-2} 让已载入的检查点适应本站自己的记录：\emph{lora} 训练小型适配器，基础
权重不动（需要 \code{peft}）；\emph{full} 更新全部 1.19 亿个参数，更容易在单个站点上过拟合。
\emph{步数}设优化器步数（300）。微调后的权重在本次会话里替换原模型；重新归一化才会用上
它们，重新载入 Chronos-2 即可恢复。Chronos-2 后端上，滚动、部分依赖和排放分解被禁用，建在
滚动之上的多尺度分解也不适用（第~\ref{sec:chronos}~节）。

% ============================ 7. NORMALISATION ================================
\chapter{归一化}\label{ch:normalise}

\begin{description}\nosep
  \item[重采样变量] 模型训练用到的每个变量，连四个日历变量一起列出。勾上的被重采样（影响
        被平均掉），没勾的保持观测值。训练完列表默认设为\emph{仅气象}——通常的去气象化；
        \emph{全部}则全勾，趋势以及季节和星期循环也一起去掉。
  \item[样本数] 蒙特卡洛抽样次数，默认 300（Chronos-2 为 8；抽多少才够，见
        第~\ref{sec:normalise}~节）。
  \item[CPU 核心数] \emph{auto} 用除一个以外的全部核心。
  \item[5--95\,\% 分位数带] 同时返回抽样的第 5、25、75 和 95 百分位，给外侧带上色。
\end{description}

序列下方，观测序列和归一化序列按一天中的小时、星期几和月份取平均，每个均值带 95\,\% 置信
区间——这就是记录的周期，和 openair 的 \code{timeVariation} 一样，也显示了天气让每个周期
变动了多少。记录显示不了的周期（逐日记录的小时）会省略；周期总是按记录本身算，与上方序列
显示的分辨率无关。脚本里写作：\code{nm.cycle\_profiles(result, ["observed",
"normalised"])}。

\screenshot{ui_normalise.png}{用全部九个气象变量、300 次抽样归一化 MY1，按日均显示。归一化
序列（红色）在三月第三周下降，那是封控，天气解释不了它。下方是它的周期：归一化序列保留了
早晚高峰和工作日与周末的差别——日历保持观测值——离散度则小于观测序列。}{fig:normalise}

图上方一行给出时刻数、抽样次数、重采样变量数，以及整条记录上归一化 $-$ 观测的均值：这里是
$-0.26$~\ugm{}，八个月下来天气的影响大致相互抵消。月均值把图里的故事讲了出来：
\begin{center}\small
\begin{tabular}{@{}lrrrrrrrr@{}}
\toprule
\NOtwo{}（\ugm{}） & 1 月 & 2 月 & 3 月 & 4 月 & 5 月 & 6 月 & 7 月 & 8 月 \\
\midrule
观测     & 60.3 & 65.2 & 51.7 & 34.5 & 30.7 & 35.1 & 31.6 & 36.6 \\
归一化   & 59.7 & 64.7 & 54.0 & 34.7 & 32.1 & 32.5 & 32.1 & 33.8 \\
\bottomrule
\end{tabular}
\end{center}
2 月到 4 月，归一化序列降了 30.0~\ugm{}（46\,\%），观测降了 30.7~\ugm{}：春季这场下降，
几乎都不是天气造成的。3 月的天气比常年干净（观测比归一化低 2.3），6 月和 8 月则更脏。

\note{用 Chronos-2 时，前 \code{context\_length} 行拿观测值填；图上用灰色阴影标出，图下方
一行说明有多少行，均值只在其余行上算。}

\caution{归一化序列是模型在典型天气下的期望。没学到天气影响的模型——拟合差，或预测变量里
根本没有气象——交回来的结果贴近观测序列，看着像一次没什么可去除的去气象化。先读\tab{模型}
标签页。}

% ============================ 8. DECOMPOSITION ================================
\chapter{分解}\label{ch:decompose}

\section{按时间尺度：排放分解}
\emph{方法}选 \emph{emission} 时，分解运行第~\ref{sec:emission}~节的嵌套归一化，画出归一化
序列及其趋势、季节、星期、日变化和短期分量，每个分量附标准差 $\sigma$（图~\ref{fig:emission}）。
图上方一行指出 $\sigma$ 最大的分量。对逐日更细的记录，日变化和星期分量画在右侧，以周期的
形式——该分量按一天中的小时和星期几取的平均，25 至 75 百分位之间加阴影——而不是几条月的
振荡曲线。

\screenshot{ui_decompose_emission.png}{MY1 的排放分解（每次归一化 100 次抽样；基准
43.1~\ugm{}）。封控同时表现为趋势（$\sigma = 9.56$）和季节分量（$\sigma = 4.74$）里的
台阶：一次性事件不是季节，而记录太短，分不开二者。日变化（右）最大（$\sigma = 9.92$）：
05:00 前低，工作时间高，离散度大是因为它在封控后减弱了，跟早晚高峰的交通一样；星期循环
（$\sigma = 3.36$）周末下降。}{fig:emission}

\section{按变量：气象分解}
\emph{方法}选 \emph{meteorology} 时，分解把观测 $-$ 归一化摊到各个气象变量上
（第~\ref{sec:metdecomp}~节）。\emph{归因方式}这时可用，\emph{分组}随之出现（选
\emph{emission} 时隐藏）：
\begin{description}\nosep
  \item[归因方式] \emph{sequential}（顺序）、\emph{shapley}，或 \emph{auto}（有分组就用
        Shapley，否则顺序）。
  \item[分组] 一张表，列出模型训练用到的气象变量及其分组：\emph{风}、\emph{温度与湿度}、
        \emph{边界层与气压}、\emph{云、辐射与降水}，也可以自己输名字。\emph{建议分组}把
        每个认得出的变量放进它的组；\emph{各自单列}清掉分组。没分组的变量单独算一项贡献。
        同组的行标同一种颜色，表下一行说明这样分组得多少项贡献。
\end{description}
图中显示 met\_total，以及每个变量或分组一幅子图，并在 \code{met\_noise}（模型没有再现的
部分）旁边给出它们 $\sigma$ 的排序（图~\ref{fig:metdecomp}）。

\screenshot{ui_decompose_met.png}{MY1 的天气影响（观测 $-$ 归一化，
$\sigma = 16.79$~\ugm{}）按精确 Shapley 值（16 个联盟）摊到四个分组：风（\code{ws}、
\code{wd}；$\sigma = 9.92$）、热力（\code{temp}、\code{RH}；5.24）、混合（\code{blh}、
\code{atmos\_pres}；4.89）和天空（\code{tcc}、\code{tp}、\code{ssrd}；3.44）。}{fig:metdecomp}

顺序归因下分解跑 $k+1$ 次归一化，Shapley 下跑 $2^k$ 次，每次\emph{样本数}次抽样，开销随
变量数涨得很快。十个变量用 Shapley 就是 1024 次归一化：请分组。气象分解和归一化不一样，
它总是重采样训练用到的每一个气象变量，不跟\tab{归一化}页上的列表走。

\section{趋势}\label{sec:trend}
分解画完的序列会带出一个问题：它在升、在降，还是平的？这个上升是真的吗？\emph{分解}一栏下
的\emph{趋势分析}方框对所选序列跑 Mann--Kendall 和 Theil--Sen 检验；记录有自相关时——MY1
的排放分量 AR(1) 残差自相关 0.94——用块自助法代替斜率的经典误差棒。勾上\emph{每次排放分解
后运行}（默认不勾），每次排放分解结束时它还会自动再跑一遍。

\emph{序列}指明检验哪一条：分解的排放分量（默认，即气象解释不了的那部分）、\emph{分离出的
长期趋势}——第~\ref{sec:emission}~节画出的那条长期分量，周期性的东西已各自分离，是更干净的
趋势输入——气象分量，或观测目标本身。\emph{段数}决定分几段：1 段（默认）就是上面的单段
检验；2～4 段改跑\emph{分段拟合}，即对所选序列运行 \code{detect\_trend\_breaks}。它在趋势
中断的地方把序列切开（如果真断了），每段各给一条 Theil--Sen 斜率及其区间；断点要过两道门槛
才成立：BIC 得改善，最终拟合相对单段的 SSE 改善还须胜过把同一搜索在块自助法重采样上重跑
一遍（要引用的 p）。哪道门槛没过，记录就按单段报告，最优候选切点及其 p 保留下来，标签页
据此说明搜索看过了哪里。把\emph{段数}从 1 调高时，仍是默认的\emph{排放分量}会自动改选长期
趋势（更干净的输入）；这个选择留着，除非你自己改。

单段检验在\tab{Trend}标签页（图~\ref{fig:trend}）画出序列和它的 Theil--Sen 直线，再给一句
评估结论，说两个 p 值该引哪个：经典的 Mann--Kendall p 值在自相关记录上偏乐观，而这条记录
恰好自相关。两段以上则画出序列、\emph{分段}的 Theil--Sen 直线和标出的断点，评估结论说明
搜索有没有找到断点（在哪、要引用的块自助法 p 是多少），或门槛把断点拒了；表格每段一行
（起、止、n、斜率、95\,\% 区间、每年），CSV 导出和命令日志条目与单段运行相同（脚本里是
\code{trend\_breaks\_call}，不是 \code{trend\_call}）（图~\ref{fig:trend_breaks}）。

\screenshot{ui_trend.png}{MY1 分解的排放分量的趋势分析（第~\ref{sec:emission}~节）：序列和
它的 Theil--Sen 直线，两个 p 值——经典的 Mann--Kendall p 值，以及因记录自相关而给出的块
自助法 p 值——图上方的评估结论说明该引哪一个。}{fig:trend}

\screenshot{ui_trend_breaks.png}{段数 2、对 MY1 分离出的长期趋势做分段拟合：搜索在 2020-03-18 19:00（2020-03-23
封控的前奏）找到断点——相对单段的 SSE 下降 88.7\,\%，搜索的块自助法 p 为 0.001，要引用的是它。第 1 段
斜率 +0.0812 每天（95\,\% [0.0761, 0.0863]），第 2 段降到 $-$0.0121 每天（95\,\% [$-$0.0134, $-$0.0108]）。
图上方的评估结论和每段一行的表格，就是\tab{趋势}标签页的内容。}{fig:trend_breaks}

\section{事件与异常}\label{sec:eventsui}
和趋势并列的另一个问题：这条记录要往哪儿去是一回事，里面\emph{发生了什么}是另一回事。
\emph{趋势}下面的\emph{事件与异常}方框给每个时刻打分，量它偏离该记录正常行为多远，再把超过
阈值的连续时段标出来——烟霾事件、泄漏事件，也有仪器发疯的日子（\code{anomaly\_scores}、
\code{detect\_events}）。

\emph{序列}沿用趋势方框的选择：分解的排放分量（天气造成不了的事件）、它分离出的长期趋势、
气象分量，或观测目标本身。\emph{方法}决定分数怎么算：\code{iqr} 是
$|x-\text{中位数}|/\mathrm{IQR}$，稳健的单变量分数，\emph{阈值 k} 是它的倍数；
\code{isolation} 是对数值及其滞后生长做的孤立森林，能抓住单变量分数看不出的形态——尖峰串、
水平跳变——它按自己的决策边界打分，阈值一行因此变灰；\code{stl\_residual} 打的是季节分解
残差的分——相对于记录自身周期的突出，而非相对于整体分布。\emph{最短长度}按记录的步长算，
短于它的时段丢掉：取 1 保留单小时尖峰，取 3 更像事件，而不是仪器闪一下。

\tab{事件}标签页画出分数条——每个时刻的分数，被标记的窗口打上阴影，有阈值时再画一条阈值线，
阴影正好是曲线越过那条线的地方——外加一张窗口表：起点、终点、步数、最高分、平均分，按最高分
排序。对有自相关的逐小时记录，$k=3$ 也会把普通天气的尾巴标出来：把阈值或最短长度调高，直到
只剩真正的事件。两个调用在同一粒种子下都是确定的，所以表里的行和条上的阴影逐小时严格对得上。

\screenshot{ui_events.png}{MY1 排放分量（第~\ref{sec:emission}~节）的\tab{事件}标签页：每个
时刻的 iqr 分数，$k=2$ 阈值画成一条线——方框默认的 $k=3$ 在这条去气象化序列上一个窗口都标
不出，所以取 2 作为第一个能标出的值——被标记的 26 个窗口打着阴影，图上方是按最高分排序的
窗口表，最高分 2.36。阴影恰好是曲线越过那条线的地方：阈值低时，标出来的正是普通天气的尾巴，
表顶是 2020-03-10 07:00--18:00 那一段，整整一天。}{fig:events}

\section{不确定度分解}\label{sec:uncertainty}
去气象化后的估计又引出一个问题：效果估计的区间里，哪部分来自气象重采样，哪部分来自模型
形式？\emph{事件与异常}下面这一栏的\emph{不确定度分解}按来源分摊来回答，不必再跑一轮蒙特
卡洛。它训练\emph{模型数}（默认 4）个成员，各用各的种子，数据是表里载入的，变量集合是训练
用的那一套；然后每个成员都对同一组\emph{样本数}（默认 300，重采样\tab{归一化}页上勾选的
变量）次气象重采样做预测——抽样只按窗口的种子播一次，第 $j$ 次抽样在所有成员里完全一样。
每个时刻的（成员 $\times$ 抽样）预测矩阵沿两个方向正交地拆方差：抽样之间的项是气象方差，
成员之间的项是模型形式方差，交互项是某个成员对某次抽样有、而集成平均下来没有的反应。三项
之和逐时刻精确等于总方差——这个恒等式由包内测试钉住，分摊不会悄悄加不回去。\emph{成员
后端}决定每个成员怎么训练：\code{lightgbm}（默认，快速随机搜索，\emph{模型}页的预算就是
试验次数）或 \code{flaml}（每个成员各跑一次 AutoML 搜索，各花各的预算）。

\tab{不确定性}标签页（图~\ref{fig:uncertainty}）把逐时刻的方差按气象、模型形式、交互叠成
一条带，带顶上的线是总方差，图上方的结论给出要紧的一句：各来源的中位数份额。MY1 在图的
设置下用三个成员、共享 50 次抽样：中位总离散里 96\,\% 来自气象重采样，交互 3\,\%，模型
形式 0\,\%——拟合出的模型对天气的作用几乎是确定的，效果拿来对比的区间是天气的，不是模型的。
表格给出各来源的中位数方差和中位数份额；份额按整条记录取中位，单个极端日没法给分摊重新
加权。本次运行按窗口实际发出的 \code{decompose\_uncertainty} 调用记进命令日志；保存的会话
把整次运行带回来，图重新画出。开销是\emph{模型数}次拟合，再加\emph{模型数}次\emph{样本数}
抽样的归一化——和一次 \code{do\_all\_unc} 的预测量一样，不多；MY1 这套设置要一两分钟，不是
一小时。

\screenshot{ui_uncertainty.png}{MY1 逐时刻离散度的来源分摊：堆叠带（气象重采样、模型形式、
交互）的中位总方差为 144.3 (\ugm{})$^{2}$，图上方的结论说明它去了哪里——96\,\% 气象重采样、
3\,\% 交互、0\,\% 模型形式（三个成员共享 50 次抽样，lightgbm）。比这个区间还小的效果，模型
形式丝毫不能把它从天气里区分出来。}{fig:uncertainty}

% ============================ 9. ROLLING ======================================
\chapter{滚动与多尺度}\label{ch:rolling}

\section{滚动归一化}\label{sec:rollingui}
\emph{窗口}（14~天）、\emph{步长}（7~天）和\emph{样本数}（100）决定第~\ref{sec:rolling}~节
的窗口；重采样哪些变量，跟\tab{归一化}页走。标签页显示每个时刻各重叠窗口的均值，两个及以上
窗口重叠的地方，附带它们之间 $\pm 1$ 个标准差（图~\ref{fig:rolling}）。导出的结果表每个窗口
一列。

\screenshot{ui_rolling.png}{MY1 的滚动归一化：33 个两周窗口，每周起一个，每个 50 次抽样。
每个窗口只去掉在它自己那两周里变化的天气。}{fig:rolling}

\section{多尺度分解}\label{sec:multiscaleui}
多尺度一栏按三种宽度各跑一次滚动归一化——默认\emph{快窗口}（14~天）、\emph{中窗口}（90~天）
和\emph{慢窗口}（365~天）——再把它们相互之间、以及与归一化序列求差
（第~\ref{sec:multiscale}~节）；所以得先跑归一化。它用\tab{滚动}页上的\emph{样本数}，开销
大约是三倍滚动运行。评估结论指出能量（$\sigma$）最大的频带；算不出来的频带会在说明里列出，
并附原因。

\screenshot{ui_multiscale.png}{MY1 在 14、60 和 120~天下的多尺度分解（默认的 365 天比这条
八个月的记录还长）：分别 77、13 和 5 个窗口。天气尺度频带承载最多（$D_\text{fast}$，
$\sigma \approx 3.0$），其次是季节内和慢频带。}{fig:multiscale}

% ============================ 10. PDP =========================================
\chapter{部分依赖}\label{ch:pdp}

训练完，变量列表里是每个预测变量和日历变量，前四个已勾好。\emph{时间变量}和\emph{仅气象}
在这两组之间切换；\emph{网格点数}（50）管曲线画多细。每幅子图画出曲线（\code{pdp\_mean}）
和 $\pm$ \code{pdp\_std} 的带，即该取值下各训练行预测值的离散度（第~\ref{sec:pdptheory}~节）。
底部短刻度是该变量在训练数据里的十分位数：刻度稀的地方——曲线两头——数据少。只取少数几个
整数值的变量（小时、星期几）画成阶梯，每个取值一个点。

\screenshot{ui_pdp.png}{MY1 模型的部分依赖：\NOtwo{} 随风速增大而下降，随边界层从最浅处
加深而下降，在约 20\,$^\circ$C 以上上升，并跟随交通的一天——05:00 前低，07:00 到 17:00
高。}{fig:pdp}

这些带宽，是因为其他预测变量在各行之间还在变；它们不是曲线的置信区间。

% ============================ 11. SCM STUDIO ==================================
\chapter{合成控制（SCM 工作室）}\label{ch:scm}

合成控制的输入不一样——是许多单元组成的长面板——所以它有自己的窗口，可以从操作栏上的
\emph{SCM 工作室}、\tab{分析 $\rightarrow$ 合成控制（SCM）工作室}（Ctrl+Shift+S）或欢迎页
上的示例卡片打开。布局和主窗口一样：左边参数面板，各页从顶部标签栏选——\tab{数据}（面板
数据及其列）、\tab{设计}、\tab{估计器}、\tab{推断}——右边结果标签页（\tab{面板}、\tab{拟合}、
\tab{拟合检查}、\tab{空间安慰剂}、\tab{时间安慰剂}、\tab{不确定性}、\tab{全部单元}、
\tab{相似度}），上方是\emph{导出当前标签页数据}。方法见第~\ref{sec:scmtheory}~节。

本章的图重现 \nm{} 论文里的 ULEZ 案例：伦敦超低排放区（ULEZ）2019 年 4 月 8 日启用，在
Marylebone Road（MY1）上用月均去气象化 \NOtwo{} 评估，拿英国其他地方的交通站点当供体。为了
让 COVID-19 不落进干预后时段，面板截到 2020 年 2 月，并且只留记录完整的交通站点：38 个月、
53 个供体。SCM 工作室的\tab{文件 $\rightarrow$ 载入示例 $\rightarrow$ 伦敦超低排放区}会设好
这个设计；脚本里写作：
\begin{lstlisting}[language=Python]
import normet as nm
scm_df = nm.datasets.load_scm()
df = scm_df.query("'2017-01-01' <= date < '2020-03-01'")
df = df[df["type"].str.contains("Traffic")]
counts = df.dropna(subset=["NO2_dw"]).groupby("code").size()
donors = [c for c in counts[counts == df["date"].nunique()].index if c != "MY1"]
df = df[df["code"].isin(donors + ["MY1"])]
df.to_csv("ulez_panel.csv", index=False)
\end{lstlisting}

\section{面板数据和列}
\emph{打开数据…}（\tab{文件 $\rightarrow$ 打开数据…}也行）读一张长表，第~\ref{sec:input}~节
里的格式都行：每个单元每个日期一行。主窗口的表带单元列时，\emph{使用主窗口数据}直接拿它来
用。\emph{日期}列按列名猜；\emph{单元}列取第一个非数值、且不同取值个数在 3 到行数四分之一
之间的列；\emph{结果变量}取名字以 \code{wn} 结尾的数值列（论文表格里的去气象化序列），没有
就取第一个数值列。三个都请自己检查。

\section{设计}
\begin{description}\nosep
  \item[处理单元] 干预到达的单元。
  \item[干预日期] 干预后时段的第一天。干预发生的那一个月算在哪边是个判断：这里取
        2019-04-01，即 ULEZ 启用的月份。
  \item[供体池] 其余每个单元，都勾着。干预可能直接到达、也可能通过溢出到达的单元请取消
        勾选——自己也受了处理的供体，会把反事实往处理单元那边拉。
  \item[按相似度排列供体]（需要 \code{foundation} 扩展）用 Chronos-2 嵌入每个单元干预前的
        记录，按它们与处理单元的动态余弦相似度排序，显示在\tab{相似度}标签页；\emph{勾选
        前 N 个}加\emph{应用}把它们设为供体池。
\end{description}
\tab{面板}标签页画出每个单元的结果变量，处理单元用红色，同时体检面板：日期 $\times$ 单元的
格子缺得超过 10\,\%，或干预前日期不足 10 个，呈琥珀色，否则绿色。

\screenshot{ui_scm_panel.png}{ULEZ 面板：MY1（红）与 53 个交通站点供体，月均去气象化
\NOtwo{}，2017 年 1 月至 2020 年 2 月。MY1 始终是其中最高的。}{fig:scmpanel}

\section{拟合}\label{sec:scmfit}
\emph{后端}选估计器（第~\ref{sec:scmtheory}~节）；\emph{机器学习后端}和\emph{时间预算}
（15~s）给 \code{mlscm} 用，\emph{抽样次数}（1000）给 \code{bayesian}。\emph{运行合成
控制}填充\tab{拟合}标签页：观测与合成、效应、权重最大的供体，以及诊断表。

对 \code{scm}、\code{abadie}、\code{did}、\code{mcnnm} 和 \code{robust}，这次拟合之后还要
再拟一次：设计不变，只把处理单元干预前的值按时间打乱——一条水平和离散度都跟处理单元一样、
却跟供体毫无关系的序列（第~\ref{sec:scmdiag}~节）。它的干预前 $R^2$ 就是表里
\code{pre\_r2, shuffled} 一行，也决定评估结论怎么读：
\begin{itemize}\nosep
  \item 打乱的序列也能达到 $R^2 \ge 0.5$，说明这个设计什么都能拟合，观测的 $R^2$ 不予
        解读。评估结论为琥珀色，给出 ATT 和两个 $R^2$，并引导你去看空间安慰剂和其他估计器。
  \item 否则评估结论看干预前的 $R^2$：0.8 及以上为绿色，0.5 起为琥珀色，更低为红色；
        并给出 ATT。
\end{itemize}
\code{mlscm} 和 \code{bayesian} 太慢，拟不了两次；它们的评估结论只看干预前的 $R^2$。
\code{scm} 拟合会多一行，说明它的权重只平衡岭回归残差。

\screenshot{ui_scm_fit.png}{MY1 的增广 SCM（\code{scm}）拟合：11 个干预后月份上
ATT $= -18.14$~\ugm{}。干预前 $R^2$ 为 0.999，可把 MY1 干预前的值打乱后仍能拟合到 0.9998
（两者都显示为 1.00），所以评估结论不解读这一拟合，而是指向安慰剂检验。}{fig:scmfit}

53 个供体、27 个干预前月份，这个设计正是增广 SCM 什么都能拟合的那一种。它权重最大的供体
占 99.5\,\%（$\text{HHI} = 0.99$），可这些权重只平衡岭回归模型的残差。

同一个设计上，各估计器的结果对不上。同一面板、同一干预日期、同一批供体：
\begin{center}\small
\begin{tabular}{@{}lrrrr@{}}
\toprule
后端 & 干预前 $R^2$ & 打乱后 $R^2$ & 干预前 RMSE & ATT（\ugm{}） \\
\midrule
\code{scm}    & 0.999  & 1.000   & 0.16  & $-18.14$ \\
\code{abadie} & $-2.94$ & $-3.33$ & 13.20 & $-6.83$ \\
\code{robust} & 0.37   & 0.16    & 5.27  & $-9.48$ \\
\code{did}    & $-0.23$ & $-0.09$ & 7.36  & $-20.67$ \\
\bottomrule
\end{tabular}
\end{center}
打乱那一列就是\tab{拟合}标签页算的基准：只有 \code{scm} 能让一条与供体无关的序列拟合得和
MY1 一样好，所以只有它的干预前 $R^2$ 不该解读。经典 SCM 根本拟合不了 MY1：它的权重非负、
和为一，合成单元就超不过最高的供体，而 MY1 高于每一个供体。稳健 SCM 也只能松松地拟合。
四个估计符号一致，大小差了三倍。这是关于设计的发现，不是软件的毛病：处理单元落在供体范围
之外，正是合成控制最棘手的情形，分析这种情形应该报告范围，而不是一个数字。

\section{推断}
\begin{description}
  \item[运行空间安慰剂] \tab{空间安慰剂}标签页把每个供体的安慰剂效应画成灰色，处理单元的
        红色，再在安慰剂的 95\,\% 区间里画出处理效应。$p \le 0.05$ 绿色，0.15 以内是标着
        “临界”的琥珀色，其余琥珀色。同时报告 Abadie RMSPE 比值检验：处理单元干预后/干预前
        的 RMSPE 比值、它在安慰剂里的排名和 p 值。
  \item[运行时间安慰剂] 每隔\emph{安慰剂间隔}步（7）设一个伪干预日期，前面至少有\emph{最短
        干预前时段}步（30）。伪干预后时段必须在真实干预日期之前结束，所以这要有较长的干预前
        时段；一个都放不下时，标签页会说明。$p \le 0.05$ 绿色；当 $p$ 已经等于伪干预日期
        个数允许的最小值时，评估结论会点明，并建议把步长调小。
  \item[运行前期趋势 + 事件研究] 干预前拟合平行吗？对干预前残差做斜率检验——斜率、p 值和
        RMPE 漂移比，这是拟合自身的平行趋势检查，\tab{拟合}标签页也报——并给出事件时间轴上
        的效应路径：干预前的点应当散在零附近，干预后的路径应当离开它们。
  \item[运行保形区间] 干预后平均效应的有限样本 95\,\% 区间，由效应块均值的子抽样零分布
        构建（Chernozhukov、W\"uthrich \& Zhu 2021）；子样本数由\emph{保形重抽样数}控制。
        $p \le 0.05$ 绿色。
  \item[运行留一供体检] 依次留出每个供体重新拟合：去掉一个供体时，其余权重和干预后效应
        挪了多少。靠单个供体撑着的结论会在这里露馅。
  \item[不确定性区间] 刀切法（依次去掉每个供体）或自助法（$B$ 次重拟合，默认 100）。评估
        结论给出干预后的平均效应，以及区间排除零的干预后日期占比；到一半及以上为绿色。
  \item[运行全部单元] 把每个单元都当处理单元、拿其余所有单元当供体拟合，再按干预后效应
        大小排序。处理单元应当冒尖：排进前十分之一为绿色。
\end{description}

\screenshot{ui_scm_placebo.png}{ULEZ 的空间安慰剂：没有哪个供体的安慰剂效应接近 MY1 的，
所以 $p = 1/54 = 0.019$，即 53 个安慰剂允许的最小值。每个安慰剂都用同一个估计器拟合，所以
哪怕干预前的拟合没有信息量，这一比较也是公平的。}{fig:scmplacebo}

\screenshot{ui_scm_bands.png}{刀切法区间：效应重算了 53 次，每次去掉一个供体。干预后平均
效应为 $-18.14$~\ugm{}，区间在 55\,\% 的干预后日期上排除零（评估结论说的是“天”；面板是
月数据）。}{fig:scmbands}

% ============================ 12. DATA ========================================
\chapter{导入数据}\label{ch:data}

\section{输入格式}\label{sec:input}
\emph{打开数据…}（Ctrl+O）能读 CSV 及其他分隔文本（\code{.csv}、
\code{.tsv}、\code{.txt}、\code{.dat}；分隔符是逗号、分号、制表符、竖线还是空格，有没有
小数逗号，编码是 UTF-8 还是 Windows 代码页——都从文件本身判断），以及它们的 gzip 压缩版、
Excel（\code{.xlsx}、\code{.xlsm}、\code{.xls}）、Parquet、Feather 和 JSON。表里要有：
\begin{itemize}\nosep
  \item 时间：一列 \code{date}、\code{datetime}、\code{time}、\code{timestamp}（不区分
        大小写，也可以是\code{日期}、\code{时间}），或一列日期加旁边一列时刻，或
        \code{year}、\code{month}、\code{day}（以及 \code{hour}）几列；
  \item 污染物，数值型；
  \item 气象及其他预测变量，数值型。
\end{itemize}
时间什么常见写法都行：ISO~8601；\code{31/01/2020 13:00}（日在前）或 \code{01/31/2020}
（月在前）——顺序从数值判断，判断不了就在文件名下说明；\code{31.01.2020}、
\code{2020年1月31日 13时}、\code{2020013113}、Unix 秒数，或 Excel 的序列日数。时间是怎么
读的，显示在文件名下方。每个时刻一行，逐时或更粗。一个文件里装着多个站点的表，应放进 SCM
工作室，或在脚本里用 \code{do\_all\_multisite} 处理。没有时间的表会被拒绝，并提示面板数据
请用 SCM 工作室。缺失值可以有：没有目标的行删掉，缺的预测变量补上
（第~\ref{sec:prepare}~节）。

\section{时钟：时区}\label{sec:clock}
文件名下方的\emph{时区}，是分析用的时钟，也是图和表显示的时钟：站点的当地时间，或 UTC
（默认）。没写偏移的时间就当作这个时钟上的时间——监测网络按当地时间记录的数据正是如此；
写了偏移的（\code{+08:00}、\code{Z}）换算到这个时钟上。这样模型学到的小时和星期几就是
站点自己的——夏令时期间，早晚高峰仍停在原来的钟点。以 UTC 取得的气象和轨迹（数据工作室、
传输工作室）在合并时换到这个时钟，于是每个小时的 ERA5 都对上它所属那个小时的测量值。这个
选择保留到下次会话；改它，打开的文件会重新读一遍。脚本里用
\code{nm.io.read\_time\_series(path, zone=\ldots)}，读法一样。

% The Data Studio's figure is taller than half a page: started at the foot of a
% page it is floated to a float page of its own, so the section opens a fresh one.
\clearpage
\section{数据工作室：空气质量监测网络与气象}\label{sec:datastudio}
\emph{获取数据}（Ctrl+D）不用离开应用，就能拼出一张可建模的表（图~\ref{fig:datastudio}）：
\screenshot{ui_data_studio.png}{在 EEA 网络、国家 \code{DE} 下搜索后的数据工作室：站点画在
地图上，选中的三个用红圈标出，对应的行在下方表格里高亮——“获取并合并”会把它们下载成一张表。
左侧是监测网络、污染物、日期范围和气象来源；\emph{来源设置}一组只在非英国网络下出现。}{fig:datastudio}

\begin{enumerate}\nosep
  \item \textbf{空气质量。}选一个监测网络。六个英国网络是 AURN（国家网络，约 210 个
        站点）、AQE、SAQN、WAQN、NI 和 LMAM（地方管理），加起来约 1500 个站点；它们既
        不要国家也不要密钥。\emph{EEA} 一行按国家（ISO-2，比如 \code{DE}）搜欧洲环境署的
        报告网络：几百到上千个站点，同样不要密钥。\emph{OpenAQ} 覆盖世界大部分地区，要一个
        免费 API 密钥——贴到\emph{来源设置}里，或设置 \code{OPENAQ\_API\_KEY}。勾污染物
        （默认 \PMtwo{}、\NOtwo{} 和 O\textsubscript{3}），然后\emph{查找站点}。
        \tab{站点}标签页按坐标把站点画在地图上，下面是表格：单击选中指针下的站点，按住鼠标
        拖出方框选中框内所有站点，⌘/Ctrl-单击加选或减选。表格跟着同步，搜索框同时收窄两者
        （可以输名称，或官方代码，如 \code{MAN3} 即 Manchester Piccadilly）；指针悬停在点上
        显示站点名。滚轮以指针为中心缩放，“+”“−”按钮逐级缩放，手形按钮（或按住鼠标中键）
        平移，\emph{地图全景}把站点重新收回视野，\emph{清空选择}（或 Esc）清空选择；双击
        站点放大到它的街道，右键打开菜单——选中视野内的站点、复制站点坐标。地图上方一行
        显示搜到多少站点、选中多少。站点下面是底图——OpenStreetMap、Esri 浅色底图或卫星
        影像，在地图上方的下拉框里选，都不要密钥，来源标在图下方——显示视野时才下载。
  \item \textbf{日期范围。}英国档案按整个日历年来，从各站点开始运行时起；选中站点后，日期
        范围对齐到它的记录。EEA 的记录按污染物标年份——它的英国站点止于 2019 年，因为英国
        脱欧后不再向该机构报告——日期范围对齐到所选站点覆盖的年份；OpenAQ 站点则带着它
        最后一次报告的日期。
  \item \textbf{气象。}Open-Meteo 的 ERA5 档案（不要密钥）或 Copernicus CDS（要
        \code{\textasciitilde/.cdsapirc}）。Open-Meteo 会加入 \code{t2m}、
        \code{d2m}、\code{rh2m}、\code{sp}、\code{tcc}、\code{tp}、\code{ssrd}、\code{ws}、
        \code{wd}、\code{u10} 和 \code{v10}——都是\emph{列}一栏认作气象的名字。它们以 UTC
        给出，合并时换到窗口的时钟上（第~\ref{sec:clock}~节）。
  \item \textbf{获取并合并}下载逐时测量值和站点坐标处的气象，按小时合并。选中多个站点时
        逐个处理，叠成一张带 \code{site} 列的表——建模步骤可以按站点分组的\emph{面板}。
        \tab{预览}和\tab{表格}显示结果；\emph{保存合并的 CSV…}把它写出去，\emph{发送到主
        窗口}在主窗口里打开它。
\end{enumerate}
英国档案就是 R 的 \code{importUKAQ()} 读的那套 openair \code{.RData} 文件，不经 R 直接
解析，数据因此和 openair 工作流一致。EEA 通路读该机构自己的站点图，为每个站点、每种污染物
下载一个公开的 Parquet 文件——一次下载就是整条记录。OpenAQ 按传感器作答；同一位置注册了
多个传感器时，优先取 \ugm{} 那个，数值于是与英国、EEA 数据同量纲（只有 ppm 传感器的位置
保留自己的单位，写在该行的 \code{unit} 里），而它提供六种常见污染物：勾了 NO 或
NO\textsubscript{x} 的，对它跳过，日志里记一行。数据工作室要联网；从它送来的表放在内存里，
命令日志会照实说，而不是写一条重复不了的 \code{read\_csv}。


\section{传输工作室：后向轨迹预测变量}\label{sec:transport}
本地气象不会告诉你空气从哪儿来。传输工作室把 HYSPLIT 后向轨迹变成能说明这一点的预测变量，
合并到表上（图~\ref{fig:transport}）。面板分两页：\emph{轨迹}处理已有的轨迹文件，
\emph{HYSPLIT} 用来自己跑轨迹。

\paragraph{轨迹文件。}\emph{添加文件…}读已有的 HYSPLIT \code{tdump} 文件（读它们不用装
HYSPLIT）。\emph{站点名称}给结果加一列 \code{site}，这样为多个受体点逐一构建的批次只填
各自站点的行。\emph{将短于此值的轨迹置空}（72~h）把因气象数据用完或空气离开区域而提前
结束的轨迹置空，否则它们会被当成短程传输。

\paragraph{源区。}可选的区域，可以是边界框（名称、经纬度的最小值和最大值），也可以是
GeoJSON 或 Shapefile 文件里的多边形（\code{geo} 扩展）。每个区域加一个停留时间特征：轨迹
在区域里待的时间比例。重叠的区域各自计入落在两者之内的端点。

\paragraph{特征。}每条轨迹一行，包含：\code{traj\_n\_endpoints} 和
\code{traj\_age\_max\_h}（质量：它回溯了多远）、\code{traj\_dist\_km}（直线距离）、
\code{traj\_pathlen\_km}、\code{traj\_speed\_kmh}、\code{traj\_inflow\_deg}（从受体点指向
起点的方位角）、\code{traj\_height\_mean} 和 \code{traj\_height\_min}、每个区域一个
\code{traj\_resid\_}\emph{区域名}，以及——\code{tdump} 里有这些数据时——沿途的降水、
混合层高度、湿度、气压和温度。

\paragraph{HYSPLIT（高级）。}本地有 \code{hyts\_std} 时，工作室可以在 \emph{HYSPLIT} 页上
自己给一张受体点表跑轨迹，可选高度、回溯小时数和时间步长，气象用它下载的 GDAS1 或磁盘上
已有的 ARL 文件。特征用\emph{轨迹}页上设的源区构建。

\paragraph{合并。}\emph{发送到主窗口}把特征合并到已打开的表上：每一行取\emph{已经到达}的
最近一条轨迹，且在一个释放间隔之内——按最近时间合并会把 06:00 的轨迹给 04:00，那可是还没
到达的空气——新的 \code{traj\_*} 列被勾为预测变量。轨迹的时间按 HYSPLIT 的写法是 UTC，
合并前先换到窗口的时钟上（第~\ref{sec:clock}~节）；传输工作室拿到的是表格时间的 UTC，从它
出发算轨迹。

\screenshot{ui_transport.png}{\nm{} 自带的两天 \code{tdump} 文件，为受体点 MY1 构建：8 个
受体时刻、9 个特征，每六小时释放一次。2020 年 1 月 1--2 日，来流方向从东（80$^\circ$）
转为西（270$^\circ$）。}{fig:transport}

\caution{每六小时释放一次的轨迹，在六个逐时行上保持不变。随机划分下，树模型会把它们当
时间戳使，测试集 $R^2$ 因而虚高（第~\ref{sec:split}~节）。考虑传输的模型，请用
\code{month\_ts} 评判。要把 \PMtwo{} 分成当地天气和传输两部分贡献，给气象分解两个分组：
\code{local} 和 \code{transport}（第~\ref{sec:metdecomp}~节）。}

\section{风向玫瑰图：污染物从哪里来}\label{sec:polar}
\tab{分析 $\rightarrow$ 风向玫瑰图…}按风向和风速画出已打开表里的一个变量，就像 openair 的
\code{polarPlot}（图~\ref{fig:polar}）：在每 10$^\circ$ 风向和每档风速的格子里取均值、
中位数、第 95 百分位数或最大值，画到罗盘上，北朝上。风速风向列按名字认（\code{ws} 和
\code{wd}），认不出就用 \code{u10} 和 \code{v10} 算，也可以手动选。脚本里写作：
\code{nm.polar\_plot(df, target="NO2", ws\_col="ws", wd\_col="wd")}。

\screenshot{ui_polar.png}{MY1 的 \NOtwo{} 均值随风向和风速的分布：南风和西南风时最高，
北风和东北风时最低。}{fig:polar}

% ============================ 13. EXPORT ======================================
\chapter{导出与重现}\label{ch:export}

\section{表格和报告}
\begin{description}
  \item[导出结果 CSV]（Ctrl+E）当前标签页背后的表：\tab{数据}上是载入的数据，\tab{模型}上
        是统计量，以及归一化序列及其分位数、分解、每个滚动窗口、多尺度序列、部分依赖值。
  \item[HTML 报告]（文件 $\rightarrow$ 导出 HTML 报告，或那个按钮；在它的菜单里选报告语言。）
        把当前标签页的结果做成一个自包含的 HTML 文件——图嵌在里面，没有脚本，没有外部
        链接——附模型和结果的摘要及其溯源信息：\nm{} 的版本、时间、平台和 Python 版本、
        主机名和用户名、数据的形状和哈希、参数面板的设置及其哈希，以及随机种子。不想让
        主机名和用户名随报告流传，分享前删掉它们。
\end{description}

\section{命令日志}\label{sec:commandlog}
每一个决定结果的步骤——打开表、训练、归一化、分解、滚动运行、多尺度运行、部分依赖——都以
做同样事情的 \nm{} 调用写进\tab{视图 $\rightarrow$ 命令日志…}，参数就是窗口实际用过的那些。
任务成功了才写条目，依据是实际发出的调用，而不是参数面板此刻显示的内容。\emph{复制}把日志
放到剪贴板，\emph{保存为脚本…}（\tab{文件 $\rightarrow$ 将会话保存为脚本…}也行）把它写成
一个能重复本次会话的 Python 文件，\emph{清空}清空日志。

\screenshot{ui_command_log.png}{本手册会话之后的命令日志，显示的是最后几条：多尺度运行和
部分依赖，写成能重复它们的调用。}{fig:commandlog}

本手册各图背后的会话存下来是这样（路径已缩短，省略了导入语句以及多尺度和部分依赖的调用）：
\begin{lstlisting}[language=Python]
# 21:09:36  Load example: London Marylebone Road NO2 (2020)
df = nm.datasets.load_my1()

# 21:10:09  Train model
df_prep, model = build_model(df, 'NO2', backend='flaml',
                             covariates=['ws', 'wd', 'temp', 'RH', 'atmos_pres', 'blh', 'tcc', 'tp', 'ssrd', 'date_unix', 'day_julian', 'weekday', 'hour'],
                             split_method='random', train_fraction=0.75,
                             model_config={'time_budget': 30, 'estimator_list': ['lgbm']},
                             seed=7654321)

# 21:10:11  Normalise
normalised = normalise(df_prep, model,
                       covariates=[...the same 13...],
                       variables_resample=['ws', 'wd', 'temp', 'RH', 'atmos_pres', 'blh', 'tcc', 'tp', 'ssrd'],
                       n_samples=300, n_cores=None,
                       return_quantiles=(0.05, 0.25, 0.75, 0.95))

# 21:10:22  Decompose
decomposition = decompose(df_prep, model, method='meteorology',
                          covariates=[...the same 13...], n_samples=100,
                          groups={'Wind': ['ws', 'wd'],
                                  'Temperature and humidity': ['temp', 'RH'],
                                  'Boundary layer and pressure': ['atmos_pres', 'blh'],
                                  'Sky and rain': ['tcc', 'tp', 'ssrd']})

# 21:10:25  Rolling normalisation
rolled = rolling(df_prep, model, covariates=[...],
                 variables_resample=[...the nine...],
                 window_days=14, rolling_every=7, n_samples=50)
\end{lstlisting}
再跑一遍，同一个模型给出的归一化和分解完全一样；模型本身可能略有不同，因为 FLAML 的搜索以
墙钟时间为界（第~\ref{sec:train}~节）。

\section{打印用图}\label{sec:print}
屏幕上的图按标签页大小画；缩到期刊 89~mm 的栏宽，文字就跟着变小。每幅图工具栏上的
\emph{打印用图…}按期刊印刷尺寸重排：标题、坐标轴标签、刻度标签、图例，每种文字都取期刊
指南给的字号，边距和刻度按这个尺寸重算，而不是简单缩放（图~\ref{fig:print}）：
\begin{description}\nosep
  \item[期刊] Elsevier、Nature、Copernicus（ACP、AMT 等）或 ACS（ES\&T），各自带着
        插图指南的要求和链接。
  \item[宽度] 期刊的栏宽（Nature 89 或 183~mm；Elsevier 90、140 或 190~mm；ACS 3.33 或
        7~in；Copernicus 至少 80~mm），或自定义宽度。只有一列子图的图从最窄的栏宽起，
        并排子图从最宽的起。
  \item[高度] 按图的需要——保持图自身比例，并加高到每个子图的数据区至少 12~mm，不超过
        期刊的最大高度；子图有固定形状（如并排的玫瑰图）会留下空白时，降到实际绘制的高度
        ——或手动设定。
  \item[子图标签] 无、a~b~c、(a)~(b)~(c)、A~B~C 或 (A)~(B)~(C)，可选粗体，可从任意
        字母起（组合图中的 c、d 两幅从 c 起）；同一列子图的字母对齐。
  \item[线宽、标记大小] 数据线的线宽；把标记加大时，时间序列会加上圆形标记。参考线保持
        原样。
  \item[图例] 位置（或直接拖）、文字大小、列数、边框。排成一行却比印刷图还宽的图例，会
        自动减少列数。
  \item[分子式] 按化学写法书写：NO$_2$、PM$_{2.5}$、SO$_4^{2-}$。
  \item[标题] 整图和各子图的标题按期刊惯例留给图注；勾上就画出来。
  \item[栅格分辨率] 300、600 或 1200~dpi，用于 TIFF 和 PNG。
\end{description}
设置下方的检查，按期刊指南读当前排版的图：宽度和高度、最小文字和最小上下标（Nature 要求
至少 5~pt、至多 7~pt，Elsevier 为 7~pt，ACS 为 8~pt）、所用字体（指南指定 Arial 或
Helvetica 时）、有没有相互重叠的刻度标签，以及子图数量（ACS 最多四个）；没通过的检查显示
为红色，设置一改，全部重查。预览可以适应窗口，也可以按实际印刷尺寸显示。在预览上双击标题、
坐标轴标签或子图字母，可以改它的文字和字号；双击刻度数字，调整字号和角度（还能给这个刻度
单独设文字）；双击图例设置位置，双击色标把它挪到子图另一侧，并以毫米设定它在版面上的宽和
高；图例和色标也可以直接拖到最易读的地方。\emph{保存…}写出期刊接受的格式——PDF、EPS 或
SVG（文字保留为可编辑文字，字体嵌入），TIFF 或 PNG 按所选分辨率——尺寸和显示的完全一致，
白底。对话框打开时，图里的文字用报告语言（上次在文件 $\rightarrow$ 导出 HTML 报告里选的
语言）。关掉后，图恢复成它在标签页里的样子；期刊、宽度、子图标签、分子式和分辨率，留给
下一幅图。投稿前请查期刊的最新指南。

\screenshot{ui_print.png}{为 Nature 双栏排版的归一化序列：宽 183~mm，子图已标字母，旁边是
Nature 指南对它的要求——唯一没过的是 4.9~pt 的下标，低于 Nature 要求的 5~pt：Nature 的
7~pt 上限已不留余地。}{fig:print}

\section{保存与重新打开会话}\label{sec:saved}
\tab{文件 $\rightarrow$ 保存结果…}（Ctrl+S）把会话写进一个文件（\code{.normet}）：表格、
参数面板的设置、模型、每一步的结果和命令日志。\tab{载入结果…}把它原样读回窗口——每个结果
交给它那一步自己的处理函数，标签页、评估结论和图都重新画出——再用一条通知说明恢复了什么。
每走完一步，窗口还会把会话另存为 \code{\textasciitilde/.normet/last\_session.normet}；
结果没保存窗口就关了或崩了，\tab{重新打开上次会话}能把它读回来。Chronos-2 检查点不写进去
（它是预训练模型，约 500~MB）：要继续，请重新载入。会话文件是 Python pickle：只打开你自己
或本窗口写出的文件。

\tab{文件 $\rightarrow$ 载入会话脚本…}运行命令日志写出的脚本（第~\ref{sec:commandlog}~节），
把结果以同样的方式带回窗口：表格、模型、结果，参数面板也按它做过的调用设好。脚本是 Python，
会照原样执行，所以窗口先问一声。

\section{方法段落}\label{sec:methods}
\tab{文件 $\rightarrow$ 方法段落…}把这次会话跑过的内容写成能放进论文方法部分的话，直接复制
即可（图~\ref{fig:methods}）：记录（长度、时间跨度、步长）、模型（后端、搜索、预测变量、
划分、随机种子）及其在留出集上的拟合，然后是跑过的每一步——归一化的抽样次数和重采样变量、
分解、滚动窗口、多尺度宽度、部分依赖的变量——后面附方法依据的参考文献。它依据每一步实际
运行时的调用写，不看参数面板此刻的状态，用界面语言。请对着你实际做的事核一遍。

\screenshot{ui_methods.png}{本手册会话的方法段落：记录、FLAML 模型及其留出集拟合，以及跑过
的每一步，并附参考文献。}{fig:methods}

\section{配置和历史}\label{sec:config}
\tab{文件 $\rightarrow$ 保存配置…}把参数面板的每一项设置写成 JSON——目标、预测变量、后端、
划分、预算、估计器、随机种子，以及每一步的设置——\tab{载入配置…}从 JSON 设置参数面板；
数据不在里面。\tab{文件 $\rightarrow$ 运行历史…}列出每次运行的时间、任务、目标、后端、行数，
外加一项细节（训练运行给留出集 $R^2$，归一化给抽样次数）。

% ============================ 14. ASSISTANT ===================================
\chapter{助手}\label{ch:assistant}

助手是一个在窗口里工作的语言模型。\nm{} 启动时它在右侧打开；\tab{视图 $\rightarrow$ 助手}
（Ctrl+Shift+A）关掉、再打开它。它读窗口里的东西，用的是 \nm{} 自己对窗口内容的解读——文字
形式：表格、模型在留出集上的拟合、归一化序列、分解；它也像按钮那样运行窗口的步骤，所以它做
的事和你自己的运行一样出现在屏幕上，也进命令日志。

对话和聊天一样：你的消息在右侧，回复按 Markdown 显示（代码加底色，表格画线），助手调用的
每个工具安静地占一行——运行中圆点是灰色，返回后变绿，被拒绝时变红并在下面写出原因——点一下
就能展开它返回的内容。Ctrl/Cmd $+$ 和 $-$（或按住该键滚动滚轮）调整文字大小，Ctrl/Cmd 0
复原；所选大小会被记住。你在底部的圆角输入框里写消息，框会随内容变高：回车发送，Shift+回车
换行；在框的第一行按上、最后一行按下，可召回此前发送过的消息，最新一条之下恢复草稿——像
shell 的历史；消息本身有多行时，上下键先移动光标。\emph{发送}、\emph{停止}和重新开始对话的
\emph{新对话}都在框内。

它最帮得上忙的地方，是选建模对象和解读结果。它用一张带\emph{应用}按钮的卡片\emph{提议}目标
和预测变量，说明排除了什么、为什么；你点“应用”之前，参数面板一点不变，而且只有你同意了，
它才会训练或归一化。它有一条规矩：只引用工具返回的数字。

\section{全自动}\label{sec:autoassistant}
\emph{密钥…}旁边的\emph{自动}让助手不必等你确认。打开后，它提议的目标和预测变量立即应用
（卡片上显示\emph{已应用}），再按窗口的顺序把分析一路做下去——数据、设定、训练、归一化、
分解——每一步结果一出来就读。模型在还有步骤没做时停下，会被告知还剩什么，接着做，最多八次；
连续两次提醒都没进展，这一轮就结束，免得做不了某一步的模型原地打转。调用模型失败或回复为空，
会重新请求，最多两次。每次提醒、每次重试，都在对话里留一行说明。随时可以点\emph{停止}结束；
这个选择保留到下次会话，按钮箭头下的\emph{全自动会做什么…}也写着这些说明。命令行智能体同样
会立即应用设定，并在它自己的循环允许的范围内把整个流程跑完。

\section{选择模型}
\begin{description}
  \item[API 模型] 命名为 \code{provider:model}：\code{anthropic}、\code{openai}、
        \code{deepseek}、\code{gemini}，或本地的 \code{lmstudio} 服务器——例如
        \code{anthropic:claude-sonnet-5-5}、\code{deepseek:deepseek-chat}、
        \code{lmstudio:local-model}；你自己的服务器（任何兼容 OpenAI 的地址）可在
        模型列表上用\emph{添加服务器…}命名。模型名旁的\emph{思考强度}随每条消息发送：
        \emph{auto} 不发送、沿用服务器默认；\code{deepseek} 提供 low、high、\emph{xhigh}、
        \emph{max}（其模板不接受 \emph{medium}），\code{openai}、\code{gemini} 和
        LM Studio 为 low、medium、high，Anthropic 为 low 到 \emph{max}；集群转发的
        模型和你自己的服务器按模型家族取——DeepSeek 的四档，其余 low、medium、high。
        它的密钥从 macOS 钥匙串（\emph{密钥…}把它存在那里）或该
        服务商的环境变量（\code{ANTHROPIC\_API\_KEY}、\code{OPENAI\_API\_KEY}、
        \code{DEEPSEEK\_API\_KEY}、\code{GEMINI\_API\_KEY}）读取，不写到别处。本地
        服务器不要密钥。
  \item[命令行代理] 登录到你自己的订阅：Claude Code、Codex、Google Antigravity 的
        \code{agy} 或 GitHub Copilot，可设它的思考强度。窗口通过 127.0.0.1 上的 MCP、
        用给这个窗口生成的令牌把工具交给代理；shell 不给，自己的文件也不给。
\end{description}
助手发给远程模型的是解读——摘要、统计量、图注——加上你的消息，不是表格本身。

% ============================ 15. SCRIPTING ===================================
\chapter{脚本：API、CLI 和 MCP}\label{ch:script}

\section{Python API}
窗口做的每件事，都是 \code{normet} 包里的一个函数。一次调用就能跑完准备、训练和归一化：
\begin{lstlisting}[language=Python]
import normet as nm

df = nm.datasets.load_my1()
met = ["ws", "wd", "temp", "RH", "atmos_pres", "blh", "tcc", "tp", "ssrd"]
predictors = met + ["date_unix", "day_julian", "weekday", "hour"]

out, model, df_prep = nm.do_all(
    df=df, target="NO2", backend="flaml",
    covariates=predictors, variables_resample=met,
    n_samples=300, n_cores=4,
)
print(nm.modStats(df_prep, model))   # training, testing and all rows
\end{lstlisting}
一步一步来：
\begin{lstlisting}[language=Python]
df_prep, model = nm.build_model(df, "NO2", backend="flaml", covariates=predictors,
                                model_config={"time_budget": 120})
norm = nm.normalise(df_prep, model, covariates=predictors,
                    variables_resample=met, n_samples=300,
                    return_quantiles=[0.025, 0.975])
auto = nm.normalise_auto(df_prep, model, covariates=predictors,
                         variables_resample=met)       # stops when converged
emi = nm.decompose(df_prep, model, covariates=predictors, method="emission")
mets = nm.decompose(df_prep, model, covariates=predictors, method="meteorology",
                    groups={"wind": ["ws", "wd"], "rest": met[2:]})
roll = nm.rolling(df_prep, model, covariates=predictors, variables_resample=met)
curves = nm.pdp(df_prep, model, var_list=["ws", "temp"])
\end{lstlisting}
更多函数，都在顶层导出：
\begin{center}\small
\begin{tabular}{@{}p{7.4cm}p{8.0cm}@{}}
\toprule
函数 & 作用 \\
\midrule
\code{do\_all\_unc}       & 在 \code{n\_models} 个种子上运行流水线，给出置信区间 \\
\code{do\_all\_multisite} & 并行的逐站点流水线 \\
\code{cv\_score}          & 前向滚动交叉验证 \\
\code{run\_scm}, \code{scm\_all} & 一次合成控制；把每个单元都当作处理单元 \\
\code{placebo\_in\_space}, \code{placebo\_in\_time} & 安慰剂检验 \\
\code{uncertainty\_bands}, \code{effect\_bands\_space} & 刀切法/自助法区间和安慰剂区间 \\
\code{scm\_diagnostics}, \code{conformal\_effect\_interval} & 拟合诊断；共形区间 \\
\code{bayesian\_scm}      & 贝叶斯 SCM（\code{bayes} 扩展） \\
\code{detect\_events}, \code{anomaly\_scores} & 归一化序列中的事件 \\
\code{normalise\_plot}, \code{decomposition\_stack}, \code{scm\_dashboard} & 窗口中的图 \\
\code{make\_run}, \code{save\_run}, \code{generate\_html\_report} & 溯源和报告 \\
\code{trend}, \code{mann\_kendall}, \code{theil\_sen} & 趋势分析及其块自助法修正（\tab{趋势}标签页，第~\ref{sec:trend}~节） \\
\code{detect\_trend\_breaks} & 分段趋势：水平或斜率在哪处断掉（BIC + 块自助法门槛） \\
\code{decompose\_uncertainty} & 把总离散分摊到气象、模型形式及二者交互（\tab{不确定性}标签页，第~\ref{sec:uncertainty}~节） \\
\code{make\_memory}       & 磁盘缓存 \\
\code{Chronos2Estimator}  & 零样本去气象化（\code{foundation} 扩展） \\
\bottomrule
\end{tabular}
\end{center}
每个耗时函数的配置也是一个数据类（\code{NormaliseConfig}、\code{DecomposeConfig}、
\code{RollingConfig}）；关键字参数覆盖它的字段。\url{https://normet.readthedocs.io} 上的指南
用实例讲每个领域，仓库 \code{notebooks/} 文件夹里的笔记本重现论文中的案例。

\section{命令行}\label{sec:cli}
装上 \code{cli} 扩展，\code{normet} 就能在 CSV 或 Parquet 文件上跑流水线：
\begin{lstlisting}[language=bash]
normet do-all my1.csv --target NO2 \
    --covariates ws,wd,temp,RH,atmos_pres,blh,tcc,tp,ssrd,date_unix,day_julian,weekday,hour \
    --resample-vars ws,wd,temp,RH,atmos_pres,blh,tcc,tp,ssrd \
    --n-samples 300 --out normalised.csv --run-path my1_run
normet report my1_run.joblib --out my1_report.html
normet decompose my1.csv --target NO2 --covariates ... --method emission --out emi.csv
normet scm ulez_panel.csv --treated MY1 --cutoff 2019-04-01 --date-col date \
    --unit-col code --outcome-col NO2_dw --backend scm --out scm.csv
normet cv my1.csv --target NO2 --covariates ... --n-splits 5 --out cv.csv
normet deweather site.csv --target PM2.5 --met-vars t2m,blh,u10,v10 --out dw.csv
normet info                     # version and the backends installed
\end{lstlisting}
每条命令都接受 \code{--config file.yaml}，键就是选项名；命令行上的选项覆盖文件里的。
\code{--run-path} 还会存一个 \code{NormetRun} 档案（\code{.joblib} 加一个
\code{.meta.json} 附属文件），\code{normet report} 可以把它转成 HTML 或 Markdown。命令行的
\code{scm} 提供 \code{scm}、\code{mlscm}、\code{abadie}、\code{did} 和 \code{mcnnm}。

\section{通过 MCP 使用 AI 助手}\label{sec:mcp}
装上 \code{mcp} 扩展，\code{normet-mcp} 就通过模型上下文协议把 \nm{} 交给 AI 助手：
\begin{lstlisting}[language=bash]
claude mcp add normet -- normet-mcp          # Claude Code
# Claude Desktop, claude_desktop_config.json:
#   {"mcpServers": {"normet": {"command": "normet-mcp"}}}
\end{lstlisting}
助手可以列出示例、检查表格、启动 \code{do-all}、\code{decompose}、\code{scm} 和 \code{cv}
任务、跟踪和取消它们、读取运行或结果、画图和写报告。每次拟合都是一条 \nm{} 命令，跑在自己的
进程和文件夹（\code{\$NORMET\_MCP\_HOME}，默认 \code{\textasciitilde/.normet/mcp}）里，所以
它的命令行就是它的记录，离了助手也能再跑。助手拿到的是 \nm{} 对表格或结果的解读，不是数据；
它只读 \code{\$NORMET\_MCP\_ROOTS}（默认是你的主文件夹）下的文件。服务器要求助手提议目标、
预测变量和供体池时讲出理由，等你确认，并且只引用工具返回的数字。

% ============================ APPENDICES ======================================
\appendix

\chapter{默认值}\label{app:defaults}
\begin{center}\small
\begin{longtable}{@{}p{4.9cm}p{4.3cm}p{5.8cm}@{}}
\toprule
设置 & 窗口 & 库 \\
\midrule
\endhead
划分方法 / 训练集比例          & random / 0.75 & random / 0.75 \\
FLAML 时间预算                 & 60 s & 90 s \\
FLAML 估计器                   & lgbm & lgbm \\
FLAML 指标                     & $R^2$ & $R^2$ \\
LightGBM 搜索次数              & 20 & --- \\
随机种子                       & 7654321 & 7654321 \\
归一化抽样次数                 & 300（Chronos-2：8） & 300（Chronos-2：8） \\
重采样变量                     & 识别出的气象 & 所有非日历预测变量 \\
分位数                         & 勾选时 5、25、75、95\,\% & 无 \\
分解抽样次数                   & 300 & 300 \\
归因方式                       & auto & 顺序（有分组时 Shapley） \\
滚动窗口 / 步长 / 抽样次数     & 14 d / 7 d / 100 & 14 d / 7 d / 300 \\
多尺度宽度                     & 14 / 90 / 365 d & --- \\
PDP 网格点数 / 范围            & 50 / 第 1--99 百分位 & 50 / 第 1--99 百分位 \\
\code{normalise\_auto}          & --- & 每批 100，P95 RSE $<$ 3\,\%，连续 3 次，$\le$ 5000 \\
Chronos-2 上下文 / 预测长度    & 2048 / 168 行 & 2048 / 168 行 \\
SCM 后端                       & scm & scm \\
时间安慰剂：最短干预前 / 间隔  & 30 / 7 步 & --- \\
不确定性区间                   & 刀切法，$B = 100$ & 刀切法，$B = 200$，80\,\% 供体 \\
\bottomrule
\end{longtable}
\end{center}

\chapter{normet 存放文件的位置}\label{app:files}
\begin{center}\small
\begin{tabular}{@{}ll@{}}
\toprule
内容 & 位置 \\
\midrule
最近打开的表（八张）           & \code{\textasciitilde/.normet/gui\_recent\_files.json} \\
运行历史                       & \code{\textasciitilde/.normet/gui\_run\_history.jsonl} \\
窗口几何、外观、语言、助手模型 & 系统的设置存储（\code{apai-sys/normet-gui}） \\
字体缓存（安装包）             & \code{\textasciitilde/.normet/matplotlib} \\
上次会话（每步之后）           & \code{\textasciitilde/.normet/last\_session.normet} \\
助手 API 密钥                  & macOS 钥匙串，服务名 \code{normet} \\
MCP 任务                       & \code{\$NORMET\_MCP\_HOME}，默认 \code{\textasciitilde/.normet/mcp} \\
Chronos-2 权重                 & Hugging Face 缓存 \\
\bottomrule
\end{tabular}
\end{center}
你不保存，就不会写出别的东西。打印用图对话框的选择、助手的文字大小和报告语言，跟其他设置
一起存在设置存储里。

\chapter{故障排除}\label{app:trouble}
\begin{description}
  \item[“The CSV must contain a 'date' column”] 把时间戳列改名为 \code{date}。面板数据
        （多个单元）请放到 SCM 工作室。
  \item[模型的评估结论为红色] 预测变量解释不了目标。检查气象是否已勾选、与目标的时间基准
        是否一致，延长预算，或加入估计器。用差的模型归一化，得到的不过是观测序列而已。
  \item[某个气象列没有被勾选] 只有第~\ref{sec:columns}~节列出的名字才认得出来；其他的请在\tab{数据}
        页和\tab{归一化}页上手动勾选。
  \item[Exact Shapley values over $k$ features need $2^k$ normalisations] 用 Shapley 分解
        超过十个变量：请给出分组，或改用顺序归因。
  \item[“……不在训练用到的变量之中”] 某个分组（来自读入的配置）里列了模型训练时没用的
        变量；提示会说清是哪一个。
  \item[“'rdata' is required but not installed”] \emph{查找站点}用 \code{rdata} 读取英国
        档案，它属于 \code{data} 扩展：\code{pip install "normet[data]"}。安装包已包含它。
  \item[多尺度提示先运行归一化] 它要与归一化序列求差。
  \item[某个多尺度频带被跳过] 它的宽度使记录上的窗口少于两个；请调小。
  \item[滚动和 PDP 被禁用] 它们需要 AutoML 模型；Chronos-2 后端没有可供改变的响应面。
  \item[Chronos-2 呈灰色] 安装 \code{normet[foundation]}（微调还需 \code{normet[finetune]}）。
  \item[Chronos-2 的评估结论为红色] 打乱天气几乎不改变它的预测：它在做自回归。检查气象列，
        或改用 AutoML 后端。
  \item[SCM：“The pre-treatment fit is no check on this design”] 一条被打乱、与供体毫无关系
        的序列在干预前拟合得和处理单元一样好：相对于干预前日期数，供体太多了。请用空间
        安慰剂和其他估计器评判效应，或者用更少、更精选的供体（相似度排序可以帮忙），或者
        更长的干预前时段。
  \item[SCM：“Donor pool too small”] 至少勾选两个供体。
  \item[SCM：没有有效的伪干预日期] 干预后时段必须能放在真实干预日期之前，且其前有
        \emph{最短干预前时段}步；调小\emph{最短干预前时段}或\emph{安慰剂间隔}，或把干预
        日期往后移。
  \item[通知要求先做某一步] 它说明缺什么，并把那一步放在按钮上；那一步完成后它会自行消失。
  \item[窗口在结果保存之前关闭了] \tab{文件 $\rightarrow$ 重新打开上次会话}可以恢复最后一步
        之后的会话。
  \item[参数面板不见了] 它被收起了：点击操作栏上的\emph{展开面板}，或按 Ctrl+B。
  \item[任务运行太久] \emph{取消}立刻释放窗口；计算在后台跑完，结果被丢弃。减少抽样次数、
        预算或分解的变量数。
  \item[任务失败] 对话框显示最后一行；完整回溯在\tab{日志}停靠窗里。
  \item[macOS：“无法打开 Normet”] 应用未签名：右键点击它，选择\emph{打开}，一次即可
        （第\ref{ch:install}章）。
\end{description}

\end{document}
